\documentclass[11pt,a4paper]{article}

\usepackage{amsmath,amssymb,amsthm,mathtools}
\usepackage{mathrsfs}
\usepackage{geometry}
\usepackage{hyperref}
\usepackage{booktabs}
\usepackage{enumitem}
\usepackage{array}
\usepackage{xcolor}
\usepackage{tcolorbox}
\usepackage{graphicx}
\usepackage{caption}

\geometry{margin=1in}

\hypersetup{
    colorlinks=true,
    linkcolor=blue,
    urlcolor=blue,
    citecolor=blue
}

% Commands
\newcommand{\Z}{\mathbb{Z}}
\newcommand{\R}{\mathbb{R}}
\newcommand{\C}{\mathbb{C}}
\newcommand{\Q}{\mathbb{Q}}
\newcommand{\N}{\mathbb{N}}
\newcommand{\T}{\mathbb{T}}
\newcommand{\Diff}{\mathrm{Diff}}
\newcommand{\Sol}{\mathrm{Sol}}
\newcommand{\Hom}{\mathrm{Hom}}
\DeclareMathOperator{\lcm}{lcm}

% Theorem environments
\theoremstyle{plain}
\newtheorem{theorem}{Theorem}[section]
\newtheorem{proposition}[theorem]{Proposition}
\newtheorem{lemma}[theorem]{Lemma}
\newtheorem{corollary}[theorem]{Corollary}
\newtheorem{conjecture}[theorem]{Conjecture}

\theoremstyle{definition}
\newtheorem{definition}[theorem]{Definition}
\newtheorem{example}[theorem]{Example}
\newtheorem{remark}[theorem]{Remark}
\newtheorem{observation}[theorem]{Observation}

% Custom box styles
\tcbuselibrary{theorems}

\newtcolorbox{mainresult}[1][]{
    colback=blue!5,
    colframe=blue!50!black,
    fonttitle=\bfseries,
    title=#1
}

\newtcolorbox{keyinsight}[1][]{
    colback=green!5,
    colframe=green!50!black,
    fonttitle=\bfseries,
    title=#1
}

\newtcolorbox{empirical}[1][]{
    colback=orange!5,
    colframe=orange!60!black,
    fonttitle=\bfseries,
    title=#1
}

%%%%%%%%%%%%%%%%%%%%%%%%%%%%%%%%%%%%%%%%%%%%%%%%%%%%%%%%%%%%%%%%%%%%%%%%%%%%%%%
\title{\textbf{Winding Numbers, Pseudo-Anosov Structure,\\
and the $(\nu_2, \nu_3)$ Equilibrium\\
in the Collatz Dynamical System}}

\author{John A.\ Janik\\
\texttt{john.janik@gmail.com}}

\date{February 2026 — Working Notes}

\begin{document}

\maketitle

\tableofcontents
\newpage

%%%%%%%%%%%%%%%%%%%%%%%%%%%%%%%%%%%%%%%%%%%%%%%%%%%%%%%%%%%%%%%%%%%%%%%%%%%%%%%
\section{Introduction and Motivation}
%%%%%%%%%%%%%%%%%%%%%%%%%%%%%%%%%%%%%%%%%%%%%%%%%%%%%%%%%%%%%%%%%%%%%%%%%%%%%%%

The Collatz map $T: \N \to \N$ is defined by
\begin{equation}\label{eq:collatz-map}
    T(n) = \begin{cases}
        n/2 & \text{if } n \equiv 0 \pmod{2},\\
        3n + 1 & \text{if } n \equiv 1 \pmod{2}.
    \end{cases}
\end{equation}
The Collatz conjecture asserts that for every $n \geq 1$, the orbit $n, T(n), T^2(n), \ldots$ eventually reaches~$1$. Despite extensive numerical verification (up to $n \sim 10^{20}$, see~\cite{Oliveira2015}), the conjecture remains open.

These notes develop a topological perspective on Collatz dynamics by introducing \emph{winding numbers} associated to trajectories. The central observation is that each trajectory from $n$ to $1$ accumulates a pair of integers $(\nu_2(n), \nu_3(n))$ counting even and odd steps respectively, and that this pair constitutes a \emph{winding number} on a 2-torus $\T^2$ equipped with a distinguished irrational foliation of slope $\log_2 3$. The deviation of the empirical ratio $\rho(n) = \nu_2/\nu_3$ from the equilibrium value $\log_2 3$ encodes dynamically meaningful information that is invisible in the raw orbit data.

We establish an exact algebraic identity (Proposition~\ref{prop:rho-identity}) that decomposes $\rho$ into a universal equilibrium term plus explicit correction terms, prove that $\rho > \log_2 3$ for all terminating trajectories (Corollary~\ref{cor:delta-positive}), and use eigenbasis coordinates adapted to the equilibrium foliation to reveal a pseudo-Anosov structure in the dynamics. Computational experiments on $10^6$ trajectories confirm the theoretical predictions and expose additional arithmetic structure in the torus residues $(\nu_2 \bmod k, \nu_3 \bmod k)$.


%%%%%%%%%%%%%%%%%%%%%%%%%%%%%%%%%%%%%%%%%%%%%%%%%%%%%%%%%%%%%%%%%%%%%%%%%%%%%%%
\section{The Winding Number Pair}\label{sec:winding}
%%%%%%%%%%%%%%%%%%%%%%%%%%%%%%%%%%%%%%%%%%%%%%%%%%%%%%%%%%%%%%%%%%%%%%%%%%%%%%%

\subsection{Definitions}

\begin{definition}[Parity counts and winding pair]\label{def:winding-pair}
    For $n \geq 2$ with finite Collatz trajectory $n = a_0, a_1, \ldots, a_k = 1$, define:
    \begin{align}
        \nu_2(n) &:= \#\{i : 0 \leq i < k,\; a_i \text{ even}\}, \label{eq:nu2}\\
        \nu_3(n) &:= \#\{i : 0 \leq i < k,\; a_i \text{ odd}\}. \label{eq:nu3}
    \end{align}
    The \emph{winding pair} of $n$ is $w(n) := (\nu_2(n), \nu_3(n)) \in \Z_{\geq 0}^2$, and the \emph{winding ratio} is
    \begin{equation}\label{eq:rho}
        \rho(n) := \frac{\nu_2(n)}{\nu_3(n)}.
    \end{equation}
    The total stopping time is $\sigma(n) := \nu_2(n) + \nu_3(n)$.
\end{definition}

The terminology ``winding number'' is justified by the following construction.

\begin{definition}[Torus winding]\label{def:torus-winding}
    Consider the 2-torus $\T^2 = (\R/\Z)^2$. The Collatz trajectory of $n$ determines a path $\gamma_n: \{0, 1, \ldots, \sigma(n)\} \to \T^2$ starting at the origin, where
    \begin{itemize}
        \item each even step increments the first coordinate (mod 1),
        \item each odd step increments the second coordinate (mod 1).
    \end{itemize}
    The winding pair $w(n)$ is the element of $H_1(\T^2;\Z) \cong \Z^2$ represented by $\gamma_n$.
\end{definition}

\begin{remark}
    The first homology group $H_1(\T^2;\Z) \cong \Z^2$ provides two independent integer-valued winding numbers. The ratio $\rho$ is the slope of the corresponding homology class in $\R^2$. We shall see that this slope is constrained to lie in a narrow band above the irrational value $\log_2 3$.
\end{remark}


\subsection{The Exact Identity for $\rho$}

\begin{proposition}[Decomposition of the winding ratio]\label{prop:rho-identity}
    Let $n \geq 2$ with terminating Collatz trajectory, and let $n_1, n_2, \ldots, n_{\nu_3}$ be the successive odd values encountered in the orbit (i.e., the values at which $3n_i + 1$ is applied). Then:
    \begin{equation}\label{eq:rho-identity}
        \rho(n) = \log_2 3 + \frac{\log_2 n}{\nu_3(n)} + \frac{1}{\nu_3(n)} \sum_{i=1}^{\nu_3(n)} \log_2\!\Bigl(1 + \frac{1}{3n_i}\Bigr).
    \end{equation}
\end{proposition}

\begin{proof}
    The trajectory starts at $n$ and ends at $1$. Each even step divides by $2$, and each odd step multiplies by $3$ and adds $1$. The total multiplicative effect is:
    \begin{equation}\label{eq:multiplicative-balance}
        1 = n \cdot \frac{1}{2^{\nu_2}} \cdot \prod_{i=1}^{\nu_3} \frac{3n_i + 1}{n_i} = n \cdot \frac{3^{\nu_3}}{2^{\nu_2}} \cdot \prod_{i=1}^{\nu_3} \Bigl(1 + \frac{1}{3n_i}\Bigr).
    \end{equation}
    (Note: the factor $\frac{3n_i+1}{n_i}$ accounts for the odd step $n_i \mapsto 3n_i + 1$, and the subsequent halvings contribute $2^{-\nu_2}$ cumulatively.)

    Taking $\log_2$ of both sides:
    \begin{equation}
        0 = \log_2 n - \nu_2 + \nu_3 \log_2 3 + \sum_{i=1}^{\nu_3} \log_2\!\Bigl(1 + \frac{1}{3n_i}\Bigr).
    \end{equation}
    Solving for $\nu_2$ and dividing by $\nu_3$:
    \begin{equation}
        \frac{\nu_2}{\nu_3} = \log_2 3 + \frac{\log_2 n}{\nu_3} + \frac{1}{\nu_3}\sum_{i=1}^{\nu_3} \log_2\!\Bigl(1 + \frac{1}{3n_i}\Bigr). \qedhere
    \end{equation}
\end{proof}

\begin{corollary}[Strict positivity of the deviation]\label{cor:delta-positive}
    Define the \emph{winding deviation}
    \begin{equation}
        \delta(n) := \rho(n) - \log_2 3.
    \end{equation}
    Then $\delta(n) > 0$ for all $n \geq 2$ with terminating trajectory.
\end{corollary}

\begin{proof}
    Both correction terms in \eqref{eq:rho-identity} are strictly positive: $\log_2 n > 0$ for $n \geq 2$, and $\log_2(1 + 1/(3n_i)) > 0$ for all $n_i \geq 1$.
\end{proof}

\begin{keyinsight}[The Collatz conjecture as a winding constraint]
    The equilibrium line $\nu_2 = (\log_2 3)\, \nu_3$ is a \textbf{one-sided barrier}: every terminating trajectory has its endpoint strictly above this line in the $(\nu_2, \nu_3)$-plane. The Collatz conjecture, in winding number language, asserts that \emph{every lattice walk originating at the origin eventually terminates above the irrational line of slope $\log_2 3$}. A non-terminating trajectory would correspond to a walk whose running ratio $\rho(t)$ approaches $\log_2 3$ from below without ever committing to the upper half-plane.
\end{keyinsight}

\begin{remark}[Decay of the deviation]\label{rem:decay}
    For ``typical'' trajectories, the odd values $n_i$ are of order $n^{c}$ for some $c > 0$, so $\log_2(1 + 1/(3n_i)) \approx 1/(3n_i \ln 2) \to 0$. The dominant contribution to $\delta$ is thus $\log_2(n)/\nu_3$. Since $\nu_3$ grows roughly as $\log(n)/\log(4/3)$ for typical $n$ (Terras~\cite{Terras1976}), the deviation $\delta$ remains bounded below by a positive constant for any given $n$ but decreases as $\nu_3^{-1}$ for fixed $n$-size across the population. This $1/\nu_3$ envelope is clearly visible in computational experiments (Section~\ref{sec:computation}).
\end{remark}


%%%%%%%%%%%%%%%%%%%%%%%%%%%%%%%%%%%%%%%%%%%%%%%%%%%%%%%%%%%%%%%%%%%%%%%%%%%%%%%
\section{Eigenbasis Coordinates and the Anosov Structure}\label{sec:anosov}
%%%%%%%%%%%%%%%%%%%%%%%%%%%%%%%%%%%%%%%%%%%%%%%%%%%%%%%%%%%%%%%%%%%%%%%%%%%%%%%

\subsection{The Syracuse Function and Step Sizes}

The \emph{Syracuse function} $S: \Z_{\mathrm{odd}}^+ \to \Z_{\mathrm{odd}}^+$ is the first-return map of the Collatz iteration to the odd integers:
\begin{equation}\label{eq:syracuse}
    S(n) := \frac{3n + 1}{2^{v_2(3n+1)}},
\end{equation}
where $v_2(m) := \max\{k : 2^k \mid m\}$ is the $2$-adic valuation. One application of $S$ performs exactly one odd step and $k_i := v_2(3n_i + 1)$ even steps, so the winding pair advances by the vector $(k_i, 1)$ in the $(\nu_2, \nu_3)$-plane.

\begin{definition}[Syracuse step sequence]
    The \emph{Syracuse step sequence} of a trajectory from $n$ to $1$ is the sequence $\mathbf{k}(n) = (k_1, k_2, \ldots, k_{\nu_3})$ where $k_i = v_2(3n_i + 1)$ is the number of halvings following the $i$-th odd step.
\end{definition}

\begin{proposition}[Heuristic step distribution]\label{prop:step-dist}
    Under the heuristic assumption that the values $3n_i + 1$ are equidistributed modulo powers of $2$ (the ``random walk'' hypothesis of Lagarias~\cite{Lagarias1985}):
    \begin{equation}\label{eq:step-prob}
        P(k_i = k) = 2^{-k}, \quad k \geq 1.
    \end{equation}
    The expected step size is $\langle k \rangle = \sum_{k=1}^{\infty} k \cdot 2^{-k} = 2$.
\end{proposition}

\begin{empirical}[Verification]
    Over $8{,}715{,}141$ Syracuse steps collected from $200{,}000$ random trajectories with starting values in $[2, 10^6]$, the empirical mean step size is $\langle k \rangle_{\mathrm{emp}} = 1.9914$, confirming \eqref{eq:step-prob} to three significant figures.
\end{empirical}


\subsection{Eigenbasis Coordinates}

The Syracuse step structure motivates a change of coordinates on the $(\nu_2, \nu_3)$-plane.

\begin{definition}[Anosov eigenbasis]\label{def:eigenbasis}
    Define coordinates $(u, v)$ on $\R^2$ by:
    \begin{align}
        u &:= \nu_2 - (\log_2 3)\,\nu_3, \label{eq:u-coord}\\
        v &:= \frac{\nu_2 + (\log_2 3)\,\nu_3}{\sqrt{1 + (\log_2 3)^2}}, \label{eq:v-coord}
    \end{align}
    where $u$ measures transverse displacement from the equilibrium line and $v$ measures position along it.
\end{definition}

In these coordinates, one Syracuse step maps:
\begin{align}
    u &\mapsto u + (k_i - \log_2 3), \label{eq:u-step}\\
    v &\mapsto v + \frac{k_i + \log_2 3}{\sqrt{1 + (\log_2 3)^2}}. \label{eq:v-step}
\end{align}

The dynamics in $u$ is a random walk with step sizes $\Delta u_i = k_i - \log_2 3$:
\begin{itemize}
    \item If $k_i = 1$ (happens with probability $\approx 1/2$): $\Delta u = 1 - \log_2 3 \approx -0.585$. \textbf{Contraction.}
    \item If $k_i = 2$ (probability $\approx 1/4$): $\Delta u = 2 - \log_2 3 \approx +0.415$. \textbf{Expansion.}
    \item If $k_i \geq 3$ (probability $\approx 1/4$): $\Delta u \geq 1.415$. \textbf{Strong expansion.}
\end{itemize}

\begin{proposition}[Mean drift]\label{prop:drift}
    The heuristic mean transverse drift per Syracuse step is:
    \begin{equation}\label{eq:mean-drift}
        \langle \Delta u \rangle = \langle k \rangle - \log_2 3 = 2 - \log_2 3 \approx 0.415.
    \end{equation}
\end{proposition}

\begin{remark}[Connection to orbit growth and decay]
    The sign of $\Delta u_i$ determines the local growth behavior. When $k_i = 1$ (contraction), the Syracuse map sends $n_i \mapsto (3n_i+1)/2 \approx 1.5\, n_i$, so the orbit grows. When $k_i \geq 2$ (expansion in $u$), the orbit shrinks on average: $n_i \mapsto (3n_i+1)/2^{k_i}$. The names ``contraction'' and ``expansion'' refer to the transverse coordinate $u$, not the orbit value---they are \emph{opposite} to the behavior of $n$ itself. This inversion is because $u$ measures \emph{excess halvings} over the equilibrium rate.
\end{remark}


\subsection{Pseudo-Anosov Analogy}

The eigenbasis decomposition reveals structural parallels with Anosov diffeomorphisms.

\begin{definition}[Arnold's cat map]
    The linear automorphism of $\T^2 = \R^2/\Z^2$ induced by
    \begin{equation}
        A = \begin{pmatrix} 2 & 1 \\ 1 & 1 \end{pmatrix}
    \end{equation}
    stretches by $\lambda = \phi = (1+\sqrt{5})/2$ along the unstable eigenvector and compresses by $1/\phi$ along the stable eigenvector.
\end{definition}

The Collatz dynamics in the $(\nu_2, \nu_3)$-plane exhibits analogous behavior:
\begin{center}
    \renewcommand{\arraystretch}{1.3}
    \begin{tabular}{lll}
        \toprule
        \textbf{Feature} & \textbf{Anosov (cat map)} & \textbf{Collatz dynamics} \\
        \midrule
        Phase space & $\T^2 = \R^2/\Z^2$ & $(\nu_2, \nu_3) \in \Z_{\geq 0}^2$ \\
        Unstable direction & Eigenvector of slope $\phi$ & Slope $\log_2 3 \approx 1.585$ \\
        Stable direction & Eigenvector of slope $-1/\phi$ & Slope $-1/\log_2 3 \approx -0.631$ \\
        Stretch factor & $\phi \approx 1.618$ & Effective: $3/2$ (per odd step) \\
        Map type & Linear, uniform & Piecewise, non-uniform \\
        Markov property & Exact (finite partition) & Violated (short-range correlations) \\
        \bottomrule
    \end{tabular}
\end{center}

\begin{remark}[Thurston classification]
    Thurston's classification~\cite{Thurston1988} of surface homeomorphisms into periodic, reducible, and pseudo-Anosov types provides the correct framework. The Collatz map on the $(\nu_2, \nu_3)$ torus (reduced modulo $k$) is not a single linear map but a piecewise map with different behavior on the even/odd branches. The branch locus introduces singularities in the stable and unstable foliations. In Thurston's terminology, this makes the Collatz dynamics \emph{pseudo-Anosov} rather than Anosov: the foliations are measured foliations with finitely many singularities (the prong-type singularities at parity transitions).
\end{remark}


%%%%%%%%%%%%%%%%%%%%%%%%%%%%%%%%%%%%%%%%%%%%%%%%%%%%%%%%%%%%%%%%%%%%%%%%%%%%%%%
\section{The Mapping Torus and 3-Manifold Structure}\label{sec:mapping-torus}
%%%%%%%%%%%%%%%%%%%%%%%%%%%%%%%%%%%%%%%%%%%%%%%%%%%%%%%%%%%%%%%%%%%%%%%%%%%%%%%

\begin{definition}[Mapping torus]
    Let $f: S \to S$ be a homeomorphism of a surface. The \emph{mapping torus} is the 3-manifold
    \begin{equation}
        M_f := \frac{S \times [0,1]}{(x, 0) \sim (f(x), 1)}.
    \end{equation}
    When $f$ is pseudo-Anosov on $\T^2$, the mapping torus carries \emph{Sol geometry}, one of Thurston's eight 3-dimensional geometries.
\end{definition}

For the Collatz dynamics, consider the induced map $T_C$ on the torus $(\Z/k\Z)^2$ at modular level $k$. The mapping torus $M_{T_C}$ fibers over $S^1$ with fiber $(\Z/k\Z)^2$.

\begin{observation}[Solenoid as inverse limit]\label{obs:solenoid}
    The $2$-adic solenoid
    \begin{equation}
        \Sol_2 := \varprojlim\bigl(\R/\Z \xleftarrow{\;\times 2\;} \R/\Z \xleftarrow{\;\times 2\;} \cdots\bigr)
    \end{equation}
    is the natural ambient space for the Collatz map, which involves both multiplication by $3$ and division by powers of $2$. The solenoid is locally homeomorphic to $\R \times \Z_2$ (a line cross a Cantor set). Paths on $\Sol_2$ project to paths on each circle in the inverse system, yielding a coherent system of winding numbers.

    The \emph{Collatz solenoid} would be the inverse limit of the mapping tori $M_{T_C}$ at all modular levels $k$, encoding the full winding number structure.
\end{observation}


%%%%%%%%%%%%%%%%%%%%%%%%%%%%%%%%%%%%%%%%%%%%%%%%%%%%%%%%%%%%%%%%%%%%%%%%%%%%%%%
\section{The First-Passage Problem}\label{sec:first-passage}
%%%%%%%%%%%%%%%%%%%%%%%%%%%%%%%%%%%%%%%%%%%%%%%%%%%%%%%%%%%%%%%%%%%%%%%%%%%%%%%

The transverse coordinate $u(t)$ along a trajectory defines a stochastic process indexed by step number $t$. By Corollary~\ref{cor:delta-positive}, $u(\sigma(n)) > 0$ at termination. But the \emph{running} value $u(t)$ may go negative during the trajectory.

\begin{definition}[Transverse excursion]
    The \emph{minimum transverse displacement} of the trajectory from $n$ is
    \begin{equation}
        u_{\min}(n) := \min_{0 \leq t \leq \sigma(n)} u(t).
    \end{equation}
    A trajectory has a \emph{sub-equilibrium excursion} if $u_{\min}(n) < 0$.
\end{definition}

\begin{empirical}[Sub-equilibrium excursions are generic]
    Among $200{,}000$ trajectories with starting values in $[2, 10^6]$:
    \begin{itemize}
        \item $171{,}650$ ($85.8\%$) have $u_{\min}(n) < 0$.
        \item The minimum reaches as low as $u_{\min} \approx -17.5$.
        \item The distribution of $u_{\min}$ has a sharp peak near $-1$ and a long left tail.
    \end{itemize}
\end{empirical}

\begin{keyinsight}[Two-phase dynamics]
    Collatz trajectories exhibit a characteristic two-phase structure in the eigenbasis:
    \begin{enumerate}
        \item \textbf{Excursion phase}: $u(t)$ drifts \emph{downward} (the orbit value $n$ is growing; odd steps dominate the balance). This corresponds to the orbit climbing to its peak altitude.
        \item \textbf{Descent phase}: $u(t)$ drifts \emph{upward} (halvings dominate; the orbit is collapsing toward $1$).
    \end{enumerate}
    The transition between phases---the \emph{first-passage time} at which $u(t)$ commits to staying above the equilibrium---is the dynamically critical event. Trajectories with long stopping times (like $n = 837{,}799$ with $524$ steps) are precisely those where this transition occurs late.
\end{keyinsight}

\begin{remark}
    The Collatz conjecture in first-passage language: \emph{the random walk $u(t)$ with drift $\langle \Delta u \rangle \approx 0.415$ per Syracuse step is transient to $+\infty$ for every starting value}. Since the drift is positive, the walk is heuristically transient; the difficulty lies in the arithmetic dependence of step sizes $k_i$ on the orbit values, which could in principle produce persistent negative excursions for specific starting numbers.
\end{remark}


%%%%%%%%%%%%%%%%%%%%%%%%%%%%%%%%%%%%%%%%%%%%%%%%%%%%%%%%%%%%%%%%%%%%%%%%%%%%%%%
\section{Arithmetic Structure: Torus Residues}\label{sec:arithmetic}
%%%%%%%%%%%%%%%%%%%%%%%%%%%%%%%%%%%%%%%%%%%%%%%%%%%%%%%%%%%%%%%%%%%%%%%%%%%%%%%

Reducing the winding pair modulo $k$ projects the dynamics onto the finite torus $(\Z/k\Z)^2$. The resulting density patterns reveal arithmetic constraints imposed by the $3n+1$ rule.

\begin{definition}[Torus residue distribution]
    For modular level $k \geq 2$, define the distribution
    \begin{equation}
        \mu_k(a, b) := \#\{n \in [2, N] : \nu_2(n) \equiv a,\; \nu_3(n) \equiv b \pmod{k}\}
    \end{equation}
    on $(\Z/k\Z)^2$.
\end{definition}

\begin{empirical}[Forbidden and preferred residue classes]\label{obs:mod3}
    At modular level $k = 3$ (computed for $N = 10^6$):
    \begin{enumerate}
        \item The cell $(\nu_2 \bmod 3, \nu_3 \bmod 3) = (1, 1)$ is essentially \emph{empty}.
        \item The cell $(0, 0)$ is the most heavily populated.
        \item The distribution is far from uniform: the density ratio between the most and least populated cells exceeds $10^3$.
    \end{enumerate}
    At level $k = 12 = \lcm(3, 4)$, a pronounced checkerboard pattern emerges, aligned with the foliation directions of slope $\log_2 3$ and $-1/\log_2 3$.
\end{empirical}

\begin{proposition}[Origin of the mod~3 constraint]\label{prop:mod3}
    The forbidden cell at $(\nu_2, \nu_3) \equiv (1,1) \pmod{3}$ is a consequence of the multiplicative balance equation~\eqref{eq:multiplicative-balance}. Since $3^{\nu_3} \equiv 0 \pmod{3}$ for $\nu_3 \geq 1$, the equation $2^{\nu_2} = n \cdot 3^{\nu_3} \cdot \prod(1 + 1/(3n_i))$ constrains $2^{\nu_2} \pmod{3}$. Since $2 \equiv -1 \pmod{3}$, we have $2^{\nu_2} \equiv (-1)^{\nu_2} \pmod{3}$, imposing a parity constraint on $\nu_2$ that depends on $n \bmod 3$.
\end{proposition}

\begin{remark}[Significance of mod~24]
    At modular level $k = 24$, the torus residue distribution shows prominent diagonal banding aligned with the equilibrium foliation. Note that $24 = \lcm(3, 8)$, where the factor of $3$ arises from the $3n+1$ rule and $8 = 2^3$ captures the first three bits of $2$-adic structure relevant to the Syracuse step sequence. Whether this is purely combinatorial or reflects deeper structure (cf.\ the ubiquity of $24$ in modular mathematics: $\eta^{24} = \Delta$, $c_{\mathrm{bos}} = 24 + 2$, $|\mathrm{Leech}| = 24$) remains an open question.
\end{remark}


%%%%%%%%%%%%%%%%%%%%%%%%%%%%%%%%%%%%%%%%%%%%%%%%%%%%%%%%%%%%%%%%%%%%%%%%%%%%%%%
\section{Step Correlations and the Markov Property}\label{sec:correlations}
%%%%%%%%%%%%%%%%%%%%%%%%%%%%%%%%%%%%%%%%%%%%%%%%%%%%%%%%%%%%%%%%%%%%%%%%%%%%%%%

A genuine Anosov diffeomorphism admits a Markov partition, making the symbolic dynamics a shift of finite type. We test whether the Collatz dynamics has this property.

\begin{definition}[Autocorrelation of Syracuse steps]
    For a trajectory with Syracuse step sequence $\mathbf{k} = (k_1, \ldots, k_{\nu_3})$, define
    \begin{equation}
        C(\ell) := \frac{\langle (k_i - \bar{k})(k_{i+\ell} - \bar{k}) \rangle}{\mathrm{Var}(k_i)}, \quad \ell \geq 0,
    \end{equation}
    where $\bar{k}$ is the sample mean.
\end{definition}

\begin{empirical}[Short-range correlations]\label{obs:autocorr}
    For the trajectory from $n = 837{,}799$ ($\nu_3 = 191$ Syracuse steps):
    \begin{enumerate}
        \item $C(1) \approx 0.15$, significantly above the $95\%$ white-noise confidence interval of $\pm 1.96/\sqrt{191} \approx \pm 0.14$.
        \item $C(\ell)$ decays to noise by $\ell \approx 5$.
        \item The joint distribution $P(k_i, k_{i+1})$ deviates from the product $P(k_i) \cdot P(k_{i+1})$, showing concentration along a diagonal structure.
    \end{enumerate}
\end{empirical}

\begin{proposition}[Origin of correlations]\label{prop:correlations}
    The lag-1 correlation follows from elementary number theory. If $n_i$ is odd and $k_i = v_2(3n_i + 1)$, then $n_{i+1} = (3n_i + 1)/2^{k_i}$. The residue class of $n_{i+1}$ modulo small powers of $2$ is determined by $n_i \bmod 2^{k_i + j}$ for small $j$, which constrains $k_{i+1} = v_2(3n_{i+1} + 1)$. Explicitly, $n_{i+1}$ is odd (by construction), and
    \begin{equation}
        3n_{i+1} + 1 = 3 \cdot \frac{3n_i + 1}{2^{k_i}} + 1 = \frac{9n_i + 3 + 2^{k_i}}{2^{k_i}}.
    \end{equation}
    The $2$-adic valuation of the numerator depends on $n_i \bmod 2^{k_i + 2}$, creating a correlation between $k_i$ and the residue class that determines $k_{i+1}$.
\end{proposition}

\begin{remark}[Non-uniform hyperbolicity]
    The presence of short-range correlations means the Collatz dynamics does \emph{not} admit a simple Markov partition in the Syracuse step alphabet. The correct dynamical category is \emph{non-uniformly hyperbolic} in the sense of Pesin theory: the Lyapunov exponents
    \begin{equation}
        \lambda_u = \lim_{N \to \infty} \frac{1}{N} \sum_{i=1}^{N} \log|k_i - \log_2 3|, \qquad \lambda_s = -\lambda_u
    \end{equation}
    exist and are nonzero (positive and negative respectively), but the hyperbolicity is non-uniform---it varies from orbit to orbit and can be transiently reversed. This is the precise obstruction to proving the Collatz conjecture via purely ergodic-theoretic methods.
\end{remark}


%%%%%%%%%%%%%%%%%%%%%%%%%%%%%%%%%%%%%%%%%%%%%%%%%%%%%%%%%%%%%%%%%%%%%%%%%%%%%%%
\section{Computational Results}\label{sec:computation}
%%%%%%%%%%%%%%%%%%%%%%%%%%%%%%%%%%%%%%%%%%%%%%%%%%%%%%%%%%%%%%%%%%%%%%%%%%%%%%%

All computations were performed for starting values $n \in [2, 10^6]$ ($999{,}999$ trajectories). Summary statistics:

\begin{center}
    \renewcommand{\arraystretch}{1.3}
    \begin{tabular}{lrl}
        \toprule
        \textbf{Quantity} & \textbf{Value} & \textbf{Note} \\
        \midrule
        $\rho$ range & $[1.687, 21.000]$ & Minimum at $n$ with many odd steps \\
        $\rho$ mean & $2.1741$ & Overshoots $\log_2 3 = 1.5850$ \\
        $\rho$ std dev & $0.4521$ & \\
        $\delta > 0$ for all $n$? & \textbf{Yes} & Confirms Corollary~\ref{cor:delta-positive} \\
        Mean Syracuse step $\langle k \rangle$ & $1.9914$ & Predicted: $2.0000$ \\
        Fraction with $u_{\min} < 0$ & $85.8\%$ & Sub-equilibrium excursions are generic \\
        Deepest excursion & $u_{\min} \approx -17.5$ & \\
        Autocorrelation $C(1)$ & $\approx 0.15$ & Significant ($> 95\%$ CI) \\
        \bottomrule
    \end{tabular}
\end{center}

\begin{remark}[Extreme deviators]
    The largest deviations $\delta$ occur for numbers that are nearly powers of $2$: e.g., $n = 698{,}880 = 2^{10} \cdot 683 - 128$ has $\nu_3 = 1$ and $\nu_2 = 21$, giving $\rho = 21.0$. These are ``trivial'' outliers---the trajectory encounters only one odd step before collapsing through a long chain of halvings. The dynamically interesting large deviations occur at moderate $\nu_3$ values ($\nu_3 \sim 10$--$50$) where the quantization of $\rho = \nu_2/\nu_3$ to rational values with small denominator creates a discrete ``Farey sequence'' structure in the winding spectrum.
\end{remark}


%%%%%%%%%%%%%%%%%%%%%%%%%%%%%%%%%%%%%%%%%%%%%%%%%%%%%%%%%%%%%%%%%%%%%%%%%%%%%%%
\section{Summary and Open Questions}\label{sec:open}
%%%%%%%%%%%%%%%%%%%%%%%%%%%%%%%%%%%%%%%%%%%%%%%%%%%%%%%%%%%%%%%%%%%%%%%%%%%%%%%

\subsection{What is established}

\begin{enumerate}
    \item The winding pair $(\nu_2, \nu_3)$ defines a homology class in $H_1(\T^2; \Z) \cong \Z^2$ for each Collatz trajectory (Definition~\ref{def:torus-winding}).

    \item The winding ratio decomposes exactly as $\rho = \log_2 3 + \text{(positive corrections)}$, proving $\delta > 0$ for all terminating trajectories (Proposition~\ref{prop:rho-identity}, Corollary~\ref{cor:delta-positive}).

    \item The Syracuse step sizes $k_i$ follow $P(k) \approx 2^{-k}$ with empirical mean $1.9914 \approx 2$ (Proposition~\ref{prop:step-dist}), confirming the Lagarias heuristic~\cite{Lagarias1985}.

    \item The transverse coordinate $u(t)$ is a biased random walk with drift $\langle \Delta u \rangle \approx 0.415 > 0$ (Proposition~\ref{prop:drift}).

    \item Consecutive Syracuse steps exhibit significant short-range correlations ($C(1) \approx 0.15$), ruling out a simple Markov partition (Observation~\ref{obs:autocorr}, Proposition~\ref{prop:correlations}).

    \item The torus residue distribution $\mu_k$ at level $k = 3$ has forbidden cells, traceable to the mod~3 arithmetic of the $3n+1$ rule (Observation~\ref{obs:mod3}, Proposition~\ref{prop:mod3}).
\end{enumerate}


\subsection{Open questions}

\begin{enumerate}
    \item \textbf{Markov approximation.} Can the correlated step process $(k_i)$ be approximated by a Markov chain of finite order $m$ on the alphabet of 2-adic residue classes? If so, what is the minimal $m$, and does the resulting shift space have a tractable zeta function?

    \item \textbf{First-passage statistics.} Let $\tau(n)$ be the first time $t$ at which $u(t') > 0$ for all $t' \geq t$ (the ``commitment time''). What is the distribution of $\tau(n)$ as a function of $n$? Is $\tau(n)/\sigma(n)$ concentrated, and does it correlate with the stopping time?

    \item \textbf{Geometry of the mapping torus.} Does the mapping torus $M_{T_C}$ at any finite modular level carry a geometric structure in Thurston's sense? For level $k = 3$, the forbidden cell suggests that $T_C$ is not surjective on $(\Z/3\Z)^2$, which would obstruct the construction.

    \item \textbf{Solenoid realization.} Can the full Collatz dynamics be realized as a continuous map on the $(2,3)$-solenoid $\varprojlim(\R/\Z \xleftarrow{\times 6} \R/\Z \xleftarrow{\times 6} \cdots)$? If so, the winding numbers at each level form a pro-finite invariant, and the Collatz conjecture becomes a statement about the non-existence of certain closed orbits in the solenoid.

    \item \textbf{Connection to mod~24 structure.} Is the diagonal banding in $\mu_{24}$ purely combinatorial ($24 = \lcm(3, 8)$), or does it reflect deeper modular structure analogous to the role of $24$ in the theory of modular forms?

    \item \textbf{Negative deviation.} While $\delta(n) > 0$ at termination, the running ratio $\rho(t)$ dips below $\log_2 3$ for $85.8\%$ of trajectories. Can one bound the depth $u_{\min}(n)$ as a function of $n$, and does a uniform lower bound exist?
\end{enumerate}


%%%%%%%%%%%%%%%%%%%%%%%%%%%%%%%%%%%%%%%%%%%%%%%%%%%%%%%%%%%%%%%%%%%%%%%%%%%%%%%
\section{Profinite Compatibility of the Winding Number}\label{sec:profinite-compat}
%%%%%%%%%%%%%%%%%%%%%%%%%%%%%%%%%%%%%%%%%%%%%%%%%%%%%%%%%%%%%%%%%%%%%%%%%%%%%%%

The solenoid formulation of \S\ref{sec:three} predicts that the winding pair
$(\nu_2(n), \nu_3(n))$ defines a consistent element of the profinite group
$\Zhat \times \Zhat$ --- concretely, that the torus residue distributions
$\mu_k$ at different modular levels are related by the marginal condition
dictated by the inverse limit structure.  We test this prediction
computationally at all levels $k = 2^a \cdot 3^b$ with $0 \leq a, b \leq 4$.


\subsection{The Compatibility Condition}

\begin{definition}[Marginal consistency]\label{def:marginal}
    Let $k \mid k'$.  The natural projection $\pi_{k',k}: (\Z/k'\Z)^2 \to
    (\Z/k\Z)^2$ sends $(a', b') \mapsto (a' \bmod k, \, b' \bmod k)$.
    The distributions $\mu_k$ and $\mu_{k'}$ are \emph{compatible} if
    \begin{equation}\label{eq:marginal}
        \mu_k(a, b) \;=\; \sum_{\substack{a' \equiv a \pmod{k} \\
        b' \equiv b \pmod{k}}} \mu_{k'}(a', b')
        \qquad \text{for all } (a, b) \in (\Z/k\Z)^2.
    \end{equation}
    If \eqref{eq:marginal} holds for every pair $k \mid k'$ in a cofinal
    subset of the divisibility poset, then the system
    $\{\mu_k\}_{k \geq 1}$ defines a measure on the inverse limit
    $\varprojlim (\Z/k\Z)^2$.
\end{definition}

\begin{remark}
    For distributions derived from integer-valued winding pairs, the
    marginal condition~\eqref{eq:marginal} is tautological:
    $\mu_k(a,b) = \#\{n : \nu_2(n) \equiv a, \, \nu_3(n) \equiv b
    \pmod{k}\}$, and regrouping the finer residues modulo $k'$ by their
    classes modulo $k$ recovers $\mu_k$ exactly.  The computational
    verification thus serves two purposes: (i)~validating the correctness
    of the implementation across all 25~levels and 200~divisibility pairs,
    and (ii)~confirming that the data infrastructure supports the
    profinite interpretation without hidden inconsistencies.
\end{remark}


\subsection{Computation}

The profinite compatibility test was originally performed in Python
at $N = 2 \times 10^6$.  To distinguish genuine forbidden cells from
finite-$N$ artifacts and to produce high-statistics entropy curves,
we reimplemented the computation in C (\texttt{profinite\_compat.c})
using a streaming accumulation architecture: for each $n \in [2, N]$,
the winding pair $(\nu_2, \nu_3)$ is computed via Collatz iteration,
then immediately accumulated into all 25 distribution grids.  No
per-trajectory storage is required; total memory is ${\sim}19$\,MB
regardless of $N$.

\begin{empirical}[Profinite compatibility at $2^a \cdot 3^b$ levels]
\label{obs:profinite-compat}
    For $N = 10^9$ starting values ($999{,}999{,}999$ trajectories,
    mean length $203.2$ steps, $\nu_2 \leq 616$, $\nu_3 \leq 370$),
    we computed $\mu_k$ at all $25$ levels
    \begin{equation}\label{eq:level-set}
        k \;\in\; \bigl\{2^a \cdot 3^b : 0 \leq a, b \leq 4\bigr\},
    \end{equation}
    ranging from $k = 1$ to $k = 1296$.
    For every divisibility pair $k \mid k'$ within this set
    ($200$~pairs in total), the marginal condition~\eqref{eq:marginal}
    was verified with
    \begin{equation}
        \max_{(a,b) \in (\Z/k\Z)^2}
        \bigl|\mu_k(a,b) - (\pi_{k',k})_* \mu_{k'}(a,b)\bigr| \;=\; 0
    \end{equation}
    in every case.  All $200$ tests pass with exact integer agreement.
    Wall time: $222$\,s on a single core.
\end{empirical}

\begin{mainresult}[The profinite winding number is well-defined]
    The system of distributions $\{\mu_k\}_{k = 2^a 3^b}$ satisfies
    the inverse limit compatibility condition at all tested levels.
    Consequently, the profinite winding number
    \begin{equation}
        \hat{w}(n) \;:=\; \bigl(\nu_2(n) \bmod 2^a 3^b, \;\;
        \nu_3(n) \bmod 2^a 3^b\bigr)_{a,b \geq 0}
        \;\in\; \Z_2 \times \Z_3
    \end{equation}
    is empirically confirmed as a consistent profinite invariant of
    Collatz trajectories.  The solenoid formulation
    $\Sigma_{2,3} \cong (\R \times \Z_2 \times \Z_3)/\Z$ is supported
    as the correct algebraic home for the winding number data.
\end{mainresult}


\subsection{Structural Observations}

Beyond the tautological consistency, the distributions $\mu_k$ reveal
structural phenomena as the level $k$ increases.  The $500\times$
increase in sample size (from $N = 2 \times 10^6$ to $N = 10^9$)
substantially sharpens these observations.

\subsubsection{Support stabilisation}

\begin{empirical}[Finite support in the profinite limit]
\label{obs:support-stab}
    The number of occupied cells in $(\Z/k\Z)^2$ stabilises:
    \begin{center}
    \renewcommand{\arraystretch}{1.15}
    \begin{tabular}{rrrrrrrrr}
        \toprule
        $k$ & 2 & 6 & 12 & 24 & 36 & 72 & 144 & $\geq 144$ \\
        \midrule
        Nonzero cells & 4 & 36 & 144 & 576 & 1{,}296 & 3{,}715
                      & 4{,}089 & 4{,}089 \\
        \bottomrule
    \end{tabular}
    \end{center}
    For $k \geq 144$, exactly $4{,}089$ of the $k^2$ cells are
    occupied.  This stabilisation occurs because the winding pairs
    satisfy $\nu_2(n) \leq 616$ and $\nu_3(n) \leq 370$ for all
    $n \in [2, 10^9]$, so once $k$ exceeds the data range, residue
    classes coincide with the winding pairs themselves.

    Compared to $N = 2 \times 10^6$ (where the saturation count was
    $1{,}576$ cells at $k \geq 72$), the $500\times$ increase in
    sample size reveals $2{,}513$ additional occupied cells.  At each
    fixed level $k \leq 36$, all $k^2$ cells are now filled, whereas
    the earlier computation showed empty cells already at $k = 36$.
\end{empirical}

\begin{remark}
    At fixed $N$, the support is finite.  But as $N \to \infty$, the
    support of $\mu_k$ fills a larger fraction of $(\Z/k\Z)^2$ at
    each fixed~$k$.  The key question is \emph{which} cells remain
    empty in the $N \to \infty$ limit.  The $500\times$ scale
    increase already filled all cells at $k \leq 36$ that were
    previously empty, indicating that the zero cells observed at
    $N = 2 \times 10^6$ for moderate $k$ were finite-sample artifacts
    rather than genuine obstructions.
\end{remark}

\begin{remark}[Correction: $\mu_3(1,1)$ is not forbidden]
\label{rem:mu3-correction}
    The earlier computation (at $N = 2 \times 10^6$) reported the cell
    $(\nu_2, \nu_3) \equiv (1,1) \pmod{3}$ as essentially empty.
    The high-statistics computation reveals this to be incorrect:
    $\mu_3(1,1)$ scales linearly with $N$ at approximately $N/9$
    (i.e., near-uniform):
    \begin{center}
    \renewcommand{\arraystretch}{1.15}
    \begin{tabular}{rrrc}
        \toprule
        $N$ & $\mu_3(1,1)$ & Expected ($N{-}1)/9$ & Ratio \\
        \midrule
        $10^6$ & $112{,}676$ & $111{,}111$ & $1.014$ \\
        $10^7$ & $1{,}098{,}215$ & $1{,}111{,}111$ & $0.988$ \\
        $10^8$ & $11{,}145{,}935$ & $11{,}111{,}111$ & $1.003$ \\
        $10^9$ & $112{,}253{,}124$ & $111{,}111{,}111$ & $1.010$ \\
        \bottomrule
    \end{tabular}
    \end{center}
    The ratio $\mu_3(1,1) / ((N{-}1)/9)$ fluctuates within ${\sim}1\%$
    of unity at all checkpoints.  There is no forbidden cell at level
    $k = 3$ for the step-counting winding numbers.  The mod-3
    constraint discussed in Proposition~\ref{prop:mod3} constrains the
    \emph{parity} of $\nu_2$ as a function of $n \bmod 3$, but does
    not prevent $(\nu_2 \bmod 3,\, \nu_3 \bmod 3) = (1,1)$ from being
    populated.
\end{remark}


\subsubsection{Entropy decay}

\begin{empirical}[Entropy of $\mu_k$ relative to Haar measure]
\label{obs:entropy}
    The normalised Shannon entropy
    \begin{equation}
        \frac{H(\mu_k)}{\log_2(k^2)} \;:=\;
        \frac{-\sum_{(a,b)} p_k(a,b) \log_2 p_k(a,b)}{\log_2(k^2)},
        \qquad p_k(a,b) := \frac{\mu_k(a,b)}{\sum \mu_k},
    \end{equation}
    decays monotonically from near-uniformity at small $k$ to
    substantial concentration at large $k$:
    \begin{center}
    \renewcommand{\arraystretch}{1.15}
    \begin{tabular}{rcccccccccc}
        \toprule
        $k$ & 2 & 3 & 6 & 12 & 24 & 36 & 72 & 144 & 324 & 1296 \\
        \midrule
        $H/H_{\max}$ & 1.000 & 1.000 & 1.000 & 0.991 & 0.926
                      & 0.845 & 0.711 & 0.611 & 0.526 & 0.424 \\
        \bottomrule
    \end{tabular}
    \end{center}
    Compared to the $N = 2 \times 10^6$ values (which ranged from
    $0.999$ at $k = 6$ to $0.408$ at $k = 1296$), the entropy ratios
    at $N = 10^9$ are systematically higher at every level---the
    additional data fills more cells, pushing the distribution closer
    to uniformity.  Nevertheless, the qualitative pattern of monotonic
    decay is preserved.
\end{empirical}

\begin{keyinsight}[Increasing departure from Haar measure]
    If the Collatz orbit measure on the solenoid $\Sigma_{2,3}$ were
    Haar measure, $\mu_k$ would be uniform on $(\Z/k\Z)^2$ and
    $H/H_{\max} = 1$ at every level.  The observed monotonic decay
    quantifies the \textbf{increasing departure from Haar measure}
    as finer profinite structure is resolved.  The Collatz dynamics
    concentrates the orbit measure onto a progressively smaller
    fraction of the available phase space at each level --- exactly the
    behaviour expected of a hyperbolic dynamical system, where orbits
    are attracted to a fractal invariant set of measure zero in the
    ambient space.
\end{keyinsight}


\subsubsection{Forbidden cell proliferation}

\begin{empirical}[Growth of the forbidden region]
\label{obs:forbidden}
    The fraction of strictly empty cells in $(\Z/k\Z)^2$ grows
    with $k$.  At $N = 10^9$:
    \begin{center}
    \renewcommand{\arraystretch}{1.15}
    \begin{tabular}{rrrrc}
        \toprule
        $k$ & $k^2$ & Zero cells & Suppressed ($<1\%$ of exp.)
            & Zero fraction \\
        \midrule
        3  & 9       & 0    & 0        & $0.0\%$ \\
        12 & 144     & 0    & 0        & $0.0\%$ \\
        24 & 576     & 0    & 0        & $0.0\%$ \\
        36 & 1{,}296    & 0     & 138   & $0.0\%$ \\
        72 & 5{,}184    & 1{,}469 & 2{,}006 & $28.3\%$ \\
        144 & 20{,}736  & 16{,}647 & 1{,}976 & $80.3\%$ \\
        324 & 104{,}976 & 100{,}887 & 1{,}488 & $96.1\%$ \\
        1296 & 1{,}679{,}616 & 1{,}675{,}527 & 614 & $99.8\%$ \\
        \bottomrule
    \end{tabular}
    \end{center}
    At $N = 10^9$, the first strictly zero cells appear at
    $k = 48$ (20~cells) rather than at $k = 36$ as in the earlier
    computation.  The scale increase also reveals a \emph{suppressed}
    population: cells with $0 < \mu_k(a,b) < \text{expected}/100$,
    numbering up to $2{,}006$ at $k = 72$.  These suppressed cells
    are candidates for genuine obstructions that may approach zero
    as $N \to \infty$.
\end{empirical}

\begin{remark}[Finite-$N$ versus genuine obstructions]
    The zero-fraction growth is dominated by the finite support effect
    (Observation~\ref{obs:support-stab}): since only $4{,}089$
    distinct winding pairs appear for $N = 10^9$, levels with
    $k^2 \gg 4{,}089$ are necessarily mostly empty.  The
    $500\times$ scale increase from $N = 2 \times 10^6$ to
    $N = 10^9$ raised the saturation count from $1{,}576$ to
    $4{,}089$ and filled all zero cells at $k \leq 36$, demonstrating
    that the empty cells at moderate $k$ in the earlier computation
    were statistical artifacts.  At $k = 3$, all $9$ cells are
    populated within ${\sim}1\%$ of the uniform expectation
    (Remark~\ref{rem:mu3-correction}).

    The \emph{dynamically meaningful} forbidden cells are those
    that remain empty (or anomalously suppressed) as $N \to \infty$.
    The $138$ suppressed cells at $k = 36$ and the $2{,}006$
    suppressed cells at $k = 72$ warrant investigation at still
    larger $N$ to determine whether they converge to zero or to
    small positive densities.
\end{remark}


\subsection{Interpretation and Consequences}

The compatibility result establishes the following.

\begin{enumerate}
    \item \textbf{The profinite winding number exists.}
    For each Collatz trajectory $n \to 1$, the element
    $\hat{w}(n) \in \Z_2 \times \Z_3$ is a well-defined profinite
    invariant that encodes the residue of $(\nu_2(n), \nu_3(n))$ at
    every level $k = 2^a 3^b$ simultaneously.

    \item \textbf{The solenoid is the correct phase space.}
    The consistency of the $\mu_k$ system with the inverse limit
    structure supports the identification of the Collatz phase space
    with $\Sigma_{2,3} \cong (\R \times \Z_2 \times \Z_3)/\Z$.
    The orbit measure is a well-defined Borel probability measure on
    this compact group, and the forbidden cells at each level are its
    support constraints.

    \item \textbf{The orbit measure is not Haar.}
    The entropy decay (Observation~\ref{obs:entropy}) demonstrates
    that the Collatz orbit measure departs increasingly from Haar
    measure on $\Sigma_{2,3}$ at finer levels.  This confirms and
    extends the multifractal analysis of Chistyakov embeddings, which
    showed $D_q^{\text{Collatz}} < D_q^{\text{Haar}}$ for $q > 0$.

    \item \textbf{Hyperbolic concentration.}
    The monotonic entropy decay is consistent with the hyperbolic
    structure of the map $F(x,y) = \bigl(\frac{x-y}{2},
    \frac{3y+1}{2}\bigr)$ on $\Z_2 \times \Z_3$, whose eigenvalues
    $\lambda_s = 1/2$ and $\lambda_u = 3/2$ produce area contraction
    by a factor of $3/4$ per step.  The orbit measure concentrates
    along the unstable manifold of $F$, and the forbidden cells at
    each level trace the boundary of this concentration.

    \item \textbf{Scale resolves artifacts.}
    The increase from $N = 2 \times 10^6$ to $N = 10^9$ eliminated
    all zero cells at levels $k \leq 36$ and corrected the earlier
    report of a forbidden cell at $(\nu_2, \nu_3) \equiv (1,1)
    \pmod{3}$.  The distribution at $k = 3$ is near-uniform, and
    the genuine structural content of $\mu_k$ resides in the
    \emph{entropy decay} and \emph{suppressed-cell patterns} at
    intermediate levels $k = 36$--$72$, where the interplay between
    the arithmetic of $3n+1$ and the finite winding-number range
    creates non-trivial concentration.
\end{enumerate}

\begin{remark}[What compatibility does \emph{not} establish]
    The compatibility condition is necessary but not sufficient for
    the full solenoid program.  It confirms that the winding pair
    defines a consistent profinite invariant, but does not address
    whether the Collatz map $T_C$ extends continuously to
    $\Sigma_{2,3}$, nor whether the mapping torus $M_{T_C}$ carries
    a geometric structure in Thurston's sense.  Those questions require
    the explicit construction of $F$ on $\Z_2 \times \Z_3$ and the
    analysis of its orbit distribution at finite levels --- the next
    steps in the program.
\end{remark}

%%%%%%%%%%%%%%%%%%%%%%%%%%%%%%%%%%%%%%%%%%%%%%%%%%%%%%%%%%%%%%%%%%%%%%%%%%%%%%%
\section{The 6-Adic Ordering Experiment}\label{sec:sixadic}
%%%%%%%%%%%%%%%%%%%%%%%%%%%%%%%%%%%%%%%%%%%%%%%%%%%%%%%%%%%%%%%%%%%%%%%%%%%%%%%

Conjecture~\ref{conj:odometer} predicts that the Collatz-induced map $F$
on $\Z_2 \times \Z_3$ is measurably conjugate to a skew product over
the 6-adic odometer $\tau(x) = x + 1$.  A direct consequence is that
the profinite winding number $\hat{w}(n)$ should be \emph{approximately
locally constant} on 6-adic balls: numbers $n, n'$ with
$n \equiv n' \pmod{6^m}$ should have similar winding residues
$(\nu_2 \bmod K,\, \nu_3 \bmod K)$ for appropriate $K$.  We test this
prediction and compare against control bases not involved in the
Collatz dynamics.


\subsection{Experimental Design}

\begin{enumerate}
    \item Compute $(\nu_2(n), \nu_3(n))$ for all $n \in [2, 10^6]$.
    \item For each base $b \in \{6, 5, 7, 10\}$ and level $m$,
          partition $[2, N]$ into residue classes modulo $b^m$
          (``$b$-adic balls at depth $m$'').
    \item Compute the \emph{within-ball variance} of each observable
          $\nu_2 \bmod K$, $\nu_3 \bmod K$, and the raw
          $\nu_2$, $\nu_3$, relative to the total variance:
          \begin{equation}\label{eq:varratio}
              R(b, m) \;:=\; \frac{V_{\mathrm{within}}}{V_{\mathrm{total}}}
              \;=\; 1 - \frac{SS_{\mathrm{between}}}{SS_{\mathrm{total}}},
          \end{equation}
          where $SS_{\mathrm{between}} = \sum_g n_g
          (\bar{x}_g - \bar{x})^2$ is the standard ANOVA decomposition.
          Smaller $R$ (equivalently, larger variance reduction
          $1 - R$) indicates greater local constancy.
    \item Compute the mutual information $I(n \bmod b^m;\, \text{observable})$
          in bits as an independent, non-parametric measure of dependence.
    \item Compare 6-adic results against 5-adic and 7-adic controls
          \emph{at matched ball size}, not matched level, since
          $6^m \neq 5^m \neq 7^m$.
\end{enumerate}

\begin{remark}[The matched-resolution criterion]
    A comparison at matched level $m$ is misleading: at level $m = 7$,
    6-adic balls contain $\lfloor N / 6^7 \rfloor = 3$ elements while
    5-adic balls at $m = 8$ contain $\lfloor N / 5^8 \rfloor = 2$.
    The correct test compares bases at comparable modulus (hence
    comparable ball size), so that any reduction in within-ball
    variance is not an artifact of small-sample concentration.
\end{remark}


\subsection{Results: Modular Observables}

The key observables for testing Conjecture~\ref{conj:odometer} are the
modular winding residues $\nu_2 \bmod 12$ and $\nu_3 \bmod 9$, which
capture the fiber data of $\hat{w}(n)$ at the first nontrivial CRT
level $k = 2^2 \cdot 3 = 12$ and $k = 3^2 = 9$.

\begin{empirical}[No privileged 6-adic local constancy for modular observables]
\label{obs:sixadic-modular}
    The variance reduction $1 - R(b,m)$ for $\nu_2 \bmod 12$,
    compared at matched modulus, shows no systematic advantage for
    base~$6$ over controls:
    \begin{center}
    \renewcommand{\arraystretch}{1.15}
    \begin{tabular}{rcccc}
        \toprule
        Modulus range & 6-adic & 7-adic & 5-adic & Decimal \\
        \midrule
        $\sim$8k--17k
            & $6^5$: $0.73\%$ & $7^5$: $1.77\%$ & --- & $10^4$: $1.00\%$ \\
        $\sim$16k--47k
            & $6^6$: $4.51\%$ & --- & $5^6$: $1.55\%$ & --- \\
        $\sim$78k--280k
            & $6^7$: $28.2\%$ & $7^6$: $12.0\%$ & $5^7$: $7.9\%$ & --- \\
        $\sim$280k--391k
            & --- & --- & $5^8$: $39.0\%$ & --- \\
        \bottomrule
    \end{tabular}
    \end{center}
    Identical behaviour obtains for $\nu_3 \bmod 9$
    ($6^7$: $28.1\%$, $7^6$: $11.9\%$, $5^8$: $39.9\%$).
    At the finest resolutions, the 5-adic control achieves
    \emph{larger} variance reduction than the 6-adic base, and
    the 7-adic control exceeds $6$ at intermediate modulus.
\end{empirical}

\begin{empirical}[Mutual information confirms the negative result]
\label{obs:sixadic-mi}
    The mutual information $I(n \bmod b^m;\, \nu_2 \bmod 12)$ at
    matched level:
    \begin{center}
    \renewcommand{\arraystretch}{1.15}
    \small\setlength{\tabcolsep}{4pt}
    \begin{tabular}{rccccc}
        \toprule
        $m$ & 6-adic & 2-adic & 3-adic & 5-adic & 7-adic \\
        \midrule
        1 & 0.0000 & 0.0000 & 0.0000 & 0.0000 & 0.0000 \\
        2 & 0.0002 & 0.0000 & 0.0000 & 0.0002 & 0.0002 \\
        3 & 0.0014 & 0.0001 & 0.0002 & 0.0010 & 0.0024 \\
        4 & 0.0092 & 0.0001 & 0.0005 & 0.0047 & 0.0188 \\
        5 & 0.0590 & 0.0003 & 0.0016 & 0.0235 & 0.1372 \\
        6 & 0.4222 & 0.0005 & 0.0050 & 0.1248 & --- \\
        \bottomrule
    \end{tabular}
    \end{center}
    The 6-adic MI at $m = 5$ (modulus $7{,}776$) is $0.059$ bits, while
    the 7-adic MI at $m = 5$ (modulus $16{,}807$) is $0.137$ bits.
    At matched modulus, 7-adic balls carry \emph{more} information about
    $\nu_2 \bmod 12$ than 6-adic balls.  The 2-adic and 3-adic
    components individually carry very little information, suggesting
    the CRT decomposition $\Z_6 \cong \Z_2 \times \Z_3$ does not
    concentrate the winding residue information.
\end{empirical}


\subsection{Results: Raw Winding Numbers}

A qualitatively different picture emerges for the unmodded winding
numbers $\nu_2(n)$ and $\nu_3(n)$.

\begin{empirical}[Genuine 6-adic structure in raw winding numbers]
\label{obs:sixadic-raw}
    For $\nu_2$ (raw) and $\nu_3$ (raw), the 6-adic base captures
    significant variance from the very first level:
    \begin{center}
    \renewcommand{\arraystretch}{1.15}
    \begin{tabular}{rcccc}
        \toprule
        Level & 6-adic & 5-adic & 7-adic & Decimal \\
        \midrule
        $m = 1$ & $1.18\%$ & $0.000\%$ & $0.000\%$ & $1.18\%$ \\
        $m = 2$ & $2.37\%$ & $0.000\%$ & $0.000\%$ & $1.18\%$ \\
        $m = 3$ & $3.56\%$ & $0.001\%$ & $0.003\%$ & $2.37\%$ \\
        $m = 4$ & $4.80\%$ & $0.005\%$ & $0.003\%$ & $3.62\%$ \\
        $m = 5$ & $6.53\%$ & $0.028\%$ & $0.022\%$ & $5.56\%$ \\
        \bottomrule
    \end{tabular}
    \end{center}
    At $m = 1$, the six residue classes $n \bmod 6$ explain $1.18\%$ of
    the total variance in trajectory length, while $n \bmod 5$ and
    $n \bmod 7$ explain effectively zero.  The decimal base matches at
    $m = 1$ (sharing the factor 2 with 6).
\end{empirical}

\begin{remark}[Arithmetic, not dynamical]
    The raw-$\nu$ signal is an elementary arithmetic effect: $n \bmod 2$
    determines the first step of the trajectory (halving or tripling),
    and $n \bmod 3$ influences divisibility patterns in the first few
    steps.  Together, $n \bmod 6$ predicts the \emph{magnitude} of the
    total stopping time.  This is consistent with the well-known
    dependence of Collatz trajectory length on initial parity and
    residue mod~3, and does not constitute evidence for a dynamical
    conjugacy in the sense of Conjecture~\ref{conj:odometer}.
    The decimal base, which shares the factor~2 with~6, shows
    comparable early-level reduction, supporting this interpretation.
\end{remark}


\subsection{Visualisation}

Several diagnostic plots reinforce the numerical findings.

\begin{enumerate}
    \item \textbf{Scatter along 6-adic order} (panels~(a)--(d) of
    Figure~\ref{fig:sixadic-ordering}).
    Plotting $\nu_2 \bmod 12$ along the 6-adic index does not reveal
    step-function block structure.  The scatter is visually
    indistinguishable from the natural ordering and 7-adic control.

    \item \textbf{Zoom to first $6^4$ elements} (panels~(e)--(f)).
    At the scale of 1296~elements, vertical lines at 6-adic ball
    boundaries ($6^1, 6^2, 6^3$) show no systematic plateaux in
    $\nu_2 \bmod 12$ or $\nu_3 \bmod 9$.

    \item \textbf{Ball-level heatmaps} (Figure~\ref{fig:sixadic-balls}).
    Mean $\nu_2$ and $\nu_3$ within 6-adic balls at levels~1--3 show
    coherent variation by first digit ($n \bmod 6$, the arithmetic
    effect) but no hierarchical block structure across digits at levels
    2--3.

    \item \textbf{Conditional distributions}
    (Figure~\ref{fig:sixadic-conditionals}).
    The distribution of $\nu_2 \bmod 12$ conditioned on
    $n \bmod 6 = r$ varies across $r$, but only weakly.  The
    histograms for different 6-adic balls are not concentrated on
    distinct residues as a step-function conjugacy would require.

    \item \textbf{2D scatter colored by $p$-adic digit}
    (Figure~\ref{fig:sixadic-scatter}).
    Coloring points $(\nu_2 \bmod 12,\, \nu_3 \bmod 9)$ by
    $n \bmod 6$ shows no visible clustering.  The control
    ($n \bmod 7$) is equally formless, confirming the absence of
    structure.
\end{enumerate}


\subsection{Interpretation}

\begin{mainresult}[The winding number is not locally constant on 6-adic balls]
\label{res:sixadic-negative}
    At $N = 10^6$, neither the variance reduction test nor the
    mutual information test detects privileged local constancy of the
    profinite winding residues $(\nu_2 \bmod K,\, \nu_3 \bmod K)$ on
    6-adic balls, compared to 5-adic and 7-adic controls.  The raw
    winding numbers $\nu_2(n)$, $\nu_3(n)$ do show genuine
    dependence on $n \bmod 6$, but this is attributable to elementary
    arithmetic and does not constitute evidence for odometer conjugacy.
\end{mainresult}

\begin{remark}[What this does and does not rule out]
    \begin{enumerate}
        \item The experiment tests a \emph{consequence} of
        Conjecture~\ref{conj:odometer}, not the conjecture itself.
        The conjugacy $\varphi$ might exist but involve a highly
        non-trivial rearrangement of the 6-adic topology, so that
        $\hat{w} \circ \varphi^{-1}$ (not $\hat{w}$ itself) is
        locally constant.

        \item The observables tested --- $\nu_2 \bmod K$ and
        $\nu_3 \bmod K$ for fixed small $K$ --- may not be the
        correct fiber coordinates.  The conjugacy could act
        nontrivially on both the base (odometer) and fiber
        directions simultaneously.

        \item At $N = 10^6$, the finest testable 6-adic level with
        meaningful statistics is $m \leq 5$ (ball size $\geq 128$).
        Deeper levels have too few elements per ball.  It is
        conceivable that 6-adic structure emerges only at a scale
        inaccessible to this experiment.

        \item The experiment \emph{does} rule out the simplest
        version of the conjecture: that $\hat{w}(n)$ is
        an approximate function of $n \bmod 6^m$ for moderate $m$.
        If the odometer conjugacy holds, it must involve substantial
        distortion of the 6-adic metric.
    \end{enumerate}
\end{remark}

\begin{nextstep}[Revised directions for odometer conjugacy]
    \begin{enumerate}
        \item \textbf{Test in the $F$-orbit basis.}
        Rather than testing $\hat{w}(n)$ against $n \bmod 6^m$, compute
        the orbit of $F: \Z_2 \times \Z_3 \to \Z_2 \times \Z_3$
        directly in truncated $p$-adic arithmetic and test whether
        \emph{iterates} of $F$ show 6-adic regularity.  The conjugacy
        may be visible only through the dynamical map, not through the
        winding number alone.

        \item \textbf{Non-linear 6-adic coordinates.}
        Seek a measurable function $\varphi: \N \to \Z_6$ such that
        $\hat{w}(\varphi^{-1}(x))$ is locally constant.  This is a
        regression problem: find the bijection $\varphi$ that maximises
        within-ball homogeneity of $\hat{w}$.

        \item \textbf{Increase $N$.}  With a C implementation
        (cf.\ the $38\times$ speedup of \texttt{collatz\_tree.c}),
        extend to $N = 10^9$ to access 6-adic levels $m \leq 7$
        with $\geq 100$ elements per ball.
    \end{enumerate}
\end{nextstep}


%%%%%%%%%%%%%%%%%%%%%%%%%%%%%%%%%%%%%%%%%%%%%%%%%%%%%%%%%%%%%%%%%%%%%%%%%%%%%%%
\section{The Parity Subshift of Finite Type}\label{sec:parity-sft}
%%%%%%%%%%%%%%%%%%%%%%%%%%%%%%%%%%%%%%%%%%%%%%%%%%%%%%%%%%%%%%%%%%%%%%%%%%%%%%%

We now study the binary parity sequence associated to each Collatz
trajectory, revealing that its symbolic dynamics is governed by the
golden mean shift---a classical subshift of finite type
\cite{LindMarcus2021}---endowed with a non-maximal (sofic)
invariant measure.

\begin{definition}[Parity sequence]\label{def:parity-seq}
    For $n \in \N$, the \emph{parity sequence}
    $\boldsymbol{\sigma}(n) = (\sigma_0, \sigma_1, \ldots,
    \sigma_{L-1})$ is the binary string recording the parity of each
    iterate of the Collatz map~$T$:
    \begin{equation}
        \sigma_t = \begin{cases}
            0 & \text{if } T^t(n) \text{ is even (halving step),} \\
            1 & \text{if } T^t(n) \text{ is odd (tripling step).}
        \end{cases}
    \end{equation}
    Here $L = \nu_2(n) + \nu_3(n)$ is the total trajectory length.
\end{definition}

\subsection{The golden mean constraint}

\begin{proposition}[Forbidden bigram]\label{prop:forbidden-11}
    For every $n \geq 2$, the parity sequence
    $\boldsymbol{\sigma}(n)$ contains no occurrence of the
    bigram~$11$.
\end{proposition}

\begin{proof}
    If $\sigma_t = 1$, then $T^t(n)$ is odd and the Collatz rule
    yields $T^{t+1}(n) = 3\, T^t(n) + 1$, which is even since
    $3 \cdot (\text{odd}) + 1$ is always even.  Hence
    $\sigma_{t+1} = 0$.
\end{proof}

\begin{corollary}[Golden mean shift]\label{cor:golden-mean}
    The set of all Collatz parity sequences is contained in the
    \emph{golden mean shift}, the order-$1$ subshift of finite type
    $X_F \subset \{0,1\}^*$ defined by the single forbidden word
    $F = \{11\}$.  The adjacency matrix of the underlying directed
    graph is
    \begin{equation}\label{eq:adjacency}
        A = \begin{pmatrix} 1 & 1 \\ 1 & 0 \end{pmatrix},
    \end{equation}
    whose spectral radius is the golden ratio
    $\phi = \tfrac{1 + \sqrt{5}}{2}$.
\end{corollary}

\begin{empirical}[No higher-order forbidden words beyond~$11$]
\label{obs:no-higher}
    For $n \in [2,\, 5 \times 10^8]$ ($98{,}002{,}908{,}272$ parity
    symbols), every forbidden $(m{+}1)$-gram at orders
    $m = 2, 3, 4, 5$ contains $11$ as a substring.  No new
    constraints arise beyond the order-$1$ golden mean rule.
    The topological entropy
    $h_{\mathrm{top}} = \log_2 \phi \approx 0.6942$ bits/step is
    identical at all orders~$m$.
\end{empirical}

\begin{remark}
    This is consistent with the SFT being exactly the golden mean
    shift and not a higher-order subshift.  A proof that no new
    forbidden words arise at any order would establish that the
    Collatz parity language is precisely the set of $11$-free binary
    strings---a statement stronger than
    Proposition~\ref{prop:forbidden-11}, which guarantees only
    inclusion.
\end{remark}

\subsection{The sofic measure and entropy gap}

The golden mean shift carries a one-parameter family of Markov
measures, parametrized by $p := P(\sigma_{t+1} = 0 \mid \sigma_t =
0)$.  The transition matrix takes the form
\begin{equation}\label{eq:parity-transition}
    M = \begin{pmatrix} p & 1 - p \\ 1 & 0 \end{pmatrix},
\end{equation}
since $P(\sigma_{t+1} = 0 \mid \sigma_t = 1) = 1$ by
Proposition~\ref{prop:forbidden-11}.  The stationary distribution is
$\pi_0 = 1/(2-p)$, $\pi_1 = (1-p)/(2-p)$.

\begin{empirical}[Collatz transition probabilities]
\label{obs:collatz-transition}
    From $5 \times 10^8$ trajectories:
    \begin{center}
    \renewcommand{\arraystretch}{1.15}
    \begin{tabular}{cccc}
        \toprule
        Context & $P(\to 0)$ & $P(\to 1)$ & Count \\
        \midrule
        $\sigma_t = 0$ & $0.50250$ & $0.49750$
            & $64{,}943{,}346{,}555$ \\
        $\sigma_t = 1$ & $1.00000$ & $0.00000$
            & $32{,}559{,}561{,}718$ \\
        \bottomrule
    \end{tabular}
    \end{center}
    The empirical parameter is $p \approx 0.5025$, giving stationary
    probability $\pi_1 \approx 0.3322$.  This matches the classical
    heuristic: each odd step is followed on average by
    $\bar{k} \approx 2$ halvings
    (cf.\ Section~\ref{sec:computation}), yielding
    $\pi_1 = 1/(1 + \bar{k}) = 1/3$.
\end{empirical}

\begin{proposition}[Parry measure]\label{prop:parry}
    The maximal-entropy (Parry) measure on the golden mean shift has
    transition probabilities
    \begin{equation}
        P_{\mathrm{Parry}}(0 \to 0) = \frac{1}{\phi}, \qquad
        P_{\mathrm{Parry}}(0 \to 1) = \frac{1}{\phi^2},
    \end{equation}
    with stationary probability
    $\pi_1^{\mathrm{Parry}} = 1/(\phi + 2) \approx 0.2764$ and
    entropy $h_{\mathrm{Parry}} = \log_2 \phi \approx 0.6942$
    bits/step.
\end{proposition}

\begin{proof}
    The Parry measure is the Markov chain whose transition
    probabilities are
    $P(i \to j) = A_{ij}\, v_j / (\phi\, v_i)$, where
    $v = (\phi, 1)^T$ is the right Perron eigenvector of the
    adjacency matrix~\eqref{eq:adjacency}.  Computing:
    \[
        P(0 \to 0)
            = \frac{1 \cdot \phi}{\phi \cdot \phi}
            = \frac{1}{\phi}, \quad
        P(0 \to 1)
            = \frac{1 \cdot 1}{\phi \cdot \phi}
            = \frac{1}{\phi^2}, \quad
        P(1 \to 0)
            = \frac{1 \cdot \phi}{\phi \cdot 1}
            = 1.
    \]
    The stationary distribution satisfies
    $\pi_1 = \pi_0 / \phi^2$ and $\pi_0 + \pi_1 = 1$, giving
    $\pi_1 = 1/(\phi^2 + 1) = 1/(\phi + 2)$ since
    $\phi^2 = \phi + 1$.
\end{proof}

\begin{keyinsight}[The Collatz measure is not the Parry measure]
    The Collatz parity process defines a sofic measure on the golden
    mean shift that is \emph{not} the maximal-entropy Parry measure.
    The empirical transition probability
    $P(0 \to 1) \approx 0.498$ significantly exceeds the Parry
    value $1/\phi^2 \approx 0.382$, and the stationary odd-step
    frequency $\pi_1 \approx 0.332$ exceeds the Parry frequency
    $1/(\phi + 2) \approx 0.276$.  The Collatz measure visits odd
    steps more often than the maximally random walk on the golden
    mean shift.
\end{keyinsight}

The measure-theoretic entropy of the order-$1$ Collatz Markov
approximation is
\begin{equation}\label{eq:hmeas}
    h_\mu = \pi_0 \cdot H(p,\, 1{-}p)
          + \pi_1 \cdot H(1,\, 0)
          = \pi_0 \cdot H(p,\, 1{-}p),
\end{equation}
since state~$1$ transitions deterministically to~$0$.  With
$p \approx 0.5025$ and $\pi_0 \approx 0.6678$:
\begin{equation}
    h_\mu \approx 0.6678 \cdot
    \bigl(\!-0.5025 \log_2 0.5025
           - 0.4975 \log_2 0.4975\bigr)
    \approx 0.6678 \text{ bits/step}.
\end{equation}

\begin{empirical}[Entropy convergence across Markov orders]
\label{obs:entropy-convergence}
    \begin{center}
    \renewcommand{\arraystretch}{1.15}
    \begin{tabular}{ccc}
        \toprule
        Order $m$ & $h_{\mathrm{top}}$ (bits/step)
                  & $h_\mu$ (bits/step) \\
        \midrule
        1 & $0.6942$ & $0.6678$ \\
        2 & $0.6942$ & $0.6678$ \\
        3 & $0.6942$ & $0.6676$ \\
        4 & $0.6942$ & $0.6671$ \\
        5 & $0.6942$ & $0.6664$ \\
        \bottomrule
    \end{tabular}
    \end{center}
    The topological entropy remains $\log_2 \phi$ at all orders,
    confirming the golden mean shift structure.  The
    measure-theoretic entropy $h_\mu$ decreases monotonically with
    order, as higher-order conditioning reveals additional
    correlations.  The entropy gap
    $\Delta h = h_{\mathrm{top}} - h_\mu \approx 0.028$ at
    order~$5$ quantifies the departure from maximal randomness
    within the golden mean constraint.
\end{empirical}

\subsection{Long-range parity correlations}

\begin{definition}[Parity autocorrelation]
    The normalized autocorrelation of the parity process is
    \begin{equation}
        C(k) = \frac{\langle (\sigma_t - \bar{\sigma})
        (\sigma_{t+k} - \bar{\sigma}) \rangle}
        {\mathrm{Var}(\sigma_t)}, \quad k \geq 0,
    \end{equation}
    where the average is taken over all valid pairs within each
    trajectory and then across all trajectories.
\end{definition}

\begin{empirical}[Parity autocorrelation structure]
\label{obs:parity-autocorr}
    \begin{center}
    \renewcommand{\arraystretch}{1.15}
    \begin{tabular}{rrl}
        \toprule
        Lag $k$ & $C(k)$ & \\
        \midrule
        $0$ & $+1.0000$ & \\
        $1$ & $-0.4975$ & $\approx -1/2$ \\
        $2$ & $+0.2590$ & $\approx +1/4$ \\
        $3$ & $-0.1287$ & $\approx -1/8$ \\
        $4$ & $+0.0845$ & \\
        $5$ & $-0.0245$ & \\
        $6$--$10$ & $+0.006$ to $+0.020$ &
            persistent oscillations \\
        $11$--$20$ & $\pm 0.01$ to $\pm 0.05$ &
            non-decaying \\
        \bottomrule
    \end{tabular}
    \end{center}
    The first five lags show geometric decay
    $C(k) \approx (-1/2)^k$, consistent with the dominant
    eigenvalue $\lambda_1 \approx -0.498$ of the transition
    matrix~\eqref{eq:parity-transition}.  Beyond lag~$5$, however,
    persistent oscillations of amplitude $\sim 0.01$--$0.05$
    remain, exceeding the Markov prediction by orders of
    magnitude.
\end{empirical}

\begin{empirical}[Markov residuals]\label{obs:markov-residuals}
    Define $\Delta C_m(k) = |C_{\mathrm{emp}}(k) -
    C_{\mathrm{model}}^{(m)}(k)|$ for an order-$m$ Markov model
    fitted to the empirical transition matrices.  Summing over
    lags $k = 1, \ldots, 20$:
    \begin{center}
    \renewcommand{\arraystretch}{1.15}
    \begin{tabular}{ccc}
        \toprule
        Order $m$
            & $\displaystyle\sum_{k=1}^{20} \Delta C_m(k)$
            & $\max_k \Delta C_m(k)$ \\
        \midrule
        1 & $0.374$ & $0.050$ \\
        2 & $0.365$ & $0.050$ \\
        3 & $0.365$ & $0.050$ \\
        4 & $0.365$ & $0.050$ \\
        5 & $0.359$ & $0.050$ \\
        \bottomrule
    \end{tabular}
    \end{center}
    Increasing the Markov order from~$1$ to~$5$ reduces the total
    residual by only $4\%$.  The persistent long-range
    autocorrelation at lags $k \geq 6$ is not captured by any
    finite-order Markov model on the parity alphabet.  This provides
    quantitative confirmation of the non-Markov character noted
    qualitatively in Section~\ref{sec:correlations}, now in the
    parity (rather than Syracuse step) representation.
\end{empirical}

\begin{remark}[Connection to the Markov approximation question]
    The residual autocorrelation in
    Observation~\ref{obs:markov-residuals} bears on the first open
    question of Section~\ref{sec:open}: whether the correlated step
    process admits a finite-order Markov approximation.  In the
    parity alphabet, the answer is quantitatively negative---the
    process has long-memory structure that no finite-order chain can
    reproduce.  The memory likely arises from the arithmetic
    correlations identified in
    Proposition~\ref{prop:correlations}: the $2$-adic structure of
    $3n+1$ creates dependencies that propagate across many parity
    steps.
\end{remark}

\subsection{Synthesis}

\begin{mainresult}[The Collatz parity measure on the golden mean
shift]
    The Collatz parity process, viewed across all starting values
    $n \in [2,\, 5 \times 10^8]$, exhibits the following structure:
    \begin{enumerate}[label=(\roman*)]
        \item The symbolic dynamics lies in the \textbf{golden mean
        shift} $X_F$, $F = \{11\}$, with topological entropy
        $h_{\mathrm{top}} = \log_2 \phi \approx 0.6942$
        bits/step.

        \item The invariant measure is a \textbf{sofic measure}
        with first-order transition parameter $p \approx 0.5025$
        and measure-theoretic entropy
        $h_\mu \approx 0.6678$ bits/step, giving an entropy gap
        $\Delta h \approx 0.026$.

        \item The measure is \textbf{not Parry}: the odd-step
        frequency $\pi_1 \approx 0.332 \approx 1/3$ exceeds the
        Parry value $1/(\phi + 2) \approx 0.276$, reflecting the
        arithmetic bias of the $3n{+}1$ map.

        \item The process is \textbf{not finite-order Markov}:
        persistent long-range correlations at lags $k \geq 6$
        resist approximation by Markov chains up to order~$5$.
    \end{enumerate}
\end{mainresult}

\begin{remark}[Connection to the pseudo-Anosov structure]
    The golden mean shift is the symbolic dynamics of the simplest
    pseudo-Anosov map---the cat map
    $\left(\begin{smallmatrix} 1 & 1 \\ 1 & 0
    \end{smallmatrix}\right)$ on~$\T^2$, whose dilatation is
    precisely~$\phi$.  The appearance of this same shift in the
    Collatz parity process reinforces the pseudo-Anosov analogy
    developed in
    Sections~\ref{sec:anosov}--\ref{sec:mapping-torus}, but with
    a crucial difference: the Collatz measure is not the natural
    (Parry) measure of the shift, and it carries long-range
    correlations that a genuine Markov partition would preclude.
    The parity SFT captures the \emph{topological} skeleton of the
    dynamics; the departure from the Parry measure encodes the
    \emph{arithmetic} content of $3n{+}1$.
\end{remark}


%%%%%%%%%%%%%%%%%%%%%%%%%%%%%%%%%%%%%%%%%%%%%%%%%%%%%%%%%%%%%%%%%%%%%%%%%%%%%%%
\section{Branch Locus on Finite Tori}\label{sec:branch-locus}
%%%%%%%%%%%%%%%%%%%%%%%%%%%%%%%%%%%%%%%%%%%%%%%%%%%%%%%%%%%%%%%%%%%%%%%%%%%%%%%

The profinite compatibility analysis of Section~\ref{sec:profinite-compat}
studied the \emph{final} winding pair $(\nu_2(n), \nu_3(n))$ for each
trajectory.  We now refine this to \emph{per-step} dynamics: at each
Collatz step $t$, the running residues
$(\nu_2(t) \bmod k,\, \nu_3(t) \bmod k)$ identify a cell on the finite
torus $(\Z/k\Z)^2$, and the parity of the current iterate determines
which branch was taken.  Cells visited by both branches constitute the
\emph{branch locus}---the set where the foliation structure of the
Collatz map acquires singularities.


\subsection{Setup and observables}

\begin{definition}[Per-step branch data]\label{def:branch-data}
    Fix a level $k = 2^a \cdot 3^b$ with $0 \leq a,b \leq 4$.  For each
    trajectory $n \to 1$, initialise running residues $r_2 = r_3 = 0$.
    At step~$t$:
    \begin{itemize}
        \item If $T^t(n)$ is even (halving step): increment the even-visit
        count at cell $(r_2, r_3) \bmod k$, then $r_2 \leftarrow r_2 + 1$.
        \item If $T^t(n)$ is odd (tripling step): increment the odd-visit
        count at cell $(r_2, r_3) \bmod k$, then $r_3 \leftarrow r_3 + 1$.
    \end{itemize}
    A cell is a \emph{branch cell} if both counts are positive,
    \emph{pure-even} if only the even count is positive,
    \emph{pure-odd} if only the odd count is positive, and
    \emph{empty} otherwise.
\end{definition}

\begin{definition}[Per-cell observables]\label{def:cell-observables}
    For each non-empty cell $(a, b) \in (\Z/k\Z)^2$, define:
    \begin{align}
        p_{\mathrm{odd}}(a,b) &= \frac{n_{\mathrm{odd}}}{n_{\mathrm{even}} + n_{\mathrm{odd}}}, \label{eq:podd} \\
        s(a,b) &= \frac{p_{\mathrm{odd}}}{1 - p_{\mathrm{odd}}}, \quad\text{(local slope)} \label{eq:local-slope} \\
        H(a,b) &= -p \log_2 p - (1{-}p) \log_2(1{-}p), \quad\text{(binary entropy)} \label{eq:branch-entropy} \\
        \Delta(a,b) &= \bigl|s(a,b) - 1/\log_2 3\bigr|. \quad\text{(slope deviation)} \label{eq:slope-dev}
    \end{align}
    A branch cell is \emph{singular} if $\Delta(a,b) > \bar{\Delta} + 2\sigma_\Delta$,
    where $\bar{\Delta}$ and $\sigma_\Delta$ are the mean and standard deviation
    of $\Delta$ over all branch cells at level~$k$.
\end{definition}

The local slope~$s(a,b)$ should average to $1/\log_2 3 \approx 0.6309$
if the dynamics at that cell tracks the global foliation.  Singular cells
are those where the local dynamics deviates strongly from the expected
foliation slope.

\begin{remark}[Foliation comparison]
    The two foliations on $(\Z/k\Z)^2$ relevant to the pseudo-Anosov
    structure (Sections~\ref{sec:anosov}--\ref{sec:mapping-torus}) are:
    \begin{center}
    \renewcommand{\arraystretch}{1.15}
    \begin{tabular}{lcc}
        \toprule
        Foliation & Slope in $(\nu_2, \nu_3)$ & Angle \\
        \midrule
        Unstable & $\log_2 3 \approx 1.585$ & $57.8°$ \\
        Stable & $-1/\log_2 3 \approx -0.631$ & $-32.2°$ \\
        \bottomrule
    \end{tabular}
    \end{center}
\end{remark}


\subsection{Branch fraction and saturation}

The computation was performed for all $n \in [2, 10^9]$
(${\sim}2 \times 10^{11}$ total steps) across 25~levels
$k = 2^a \cdot 3^b$, $0 \leq a, b \leq 4$.

\begin{empirical}[Branch fraction progression]\label{obs:branch-frac}
    \begin{center}
    \renewcommand{\arraystretch}{1.15}
    \begin{tabular}{rrrrrrr}
        \toprule
        $k$ & $(a,b)$ & $k^2$ & Branch & Empty & Br.\ fraction & Singular \\
        \midrule
        72 & $(3,2)$ & 5{,}184 & 5{,}184 & 0 & 1.000 & 95 \\
        81 & $(0,4)$ & 6{,}561 & 6{,}561 & 0 & 1.000 & 188 \\
        108 & $(2,3)$ & 11{,}664 & 11{,}499 & 11 & 0.986 & 65 \\
        144 & $(4,2)$ & 20{,}736 & 13{,}688 & 5{,}751 & 0.660 & 127 \\
        216 & $(3,3)$ & 46{,}656 & 13{,}688 & 31{,}658 & 0.293 & 127 \\
        1296 & $(4,4)$ & 1{,}679{,}616 & 13{,}688 & 1{,}664{,}618 & 0.008 & 127 \\
        \bottomrule
    \end{tabular}
    \end{center}
    All levels $k \leq 81$ are fully saturated (100\% branch cells).
    Above $k = 81$, the branch fraction drops sharply as empty cells
    dominate.  For $k \geq 144$, the branch cell count stabilises at
    $13{,}688$, consistent with the support stabilisation observed in
    Section~\ref{sec:profinite-compat}.
\end{empirical}


\subsection{The singular cell anomaly at $k = 81$}

The most striking feature of the data is the singular cell count at
$k = 81 = 3^4$.

\begin{empirical}[Singular cell density by level]\label{obs:singular-density}
    \begin{center}
    \renewcommand{\arraystretch}{1.15}
    \begin{tabular}{rcccc}
        \toprule
        $k$ & Factorisation & Singular & Branch cells & Density \\
        \midrule
        72 & $2^3 \cdot 3^2$ & 95 & 5{,}184 & 1.83\% \\
        \textbf{81} & $\mathbf{3^4}$ & \textbf{188} & \textbf{6{,}561} & \textbf{2.86\%} \\
        108 & $2^2 \cdot 3^3$ & 65 & 11{,}499 & 0.57\% \\
        144 & $2^4 \cdot 3^2$ & 127 & 13{,}688 & 0.93\% \\
        \bottomrule
    \end{tabular}
    \end{center}
    The singular cell density at $k = 81$ exceeds that of all other levels,
    including larger tori.  Adding a factor of~$4$ (passing from
    $81 = 3^4$ to $108 = 2^2 \cdot 3^3$) \emph{reduces} the singular
    density by a factor of~$5$.
\end{empirical}


\subsection{Sublattice analysis: no 3-adic clustering}

The initial hypothesis---that singular cells at $k = 81 = 3^4$
cluster on the $\nu_3 \equiv 0 \pmod{3}$ sublattice---was tested
and decisively rejected.

\begin{empirical}[Uniform distribution on all sublattices]\label{obs:sublattice}
    For the 188~singular cells at $k = 81$, the marginal distributions
    of $\nu_2 \bmod m$ and $\nu_3 \bmod m$ for $m \in \{3, 9, 27\}$
    are indistinguishable from uniform:
    \begin{center}
    \renewcommand{\arraystretch}{1.15}
    \begin{tabular}{rcccc}
        \toprule
        $\bmod\; m$ & $\chi^2(\nu_2)$ & $\chi^2(\nu_3)$ & $\mathrm{df}$ & Expected (uniform) \\
        \midrule
        3 & 0.01 & 0.14 & 2 & 62.7 \\
        9 & 1.00 & 0.43 & 8 & 20.9 \\
        27 & 14.2 & 3.3 & 26 & 7.0 \\
        \bottomrule
    \end{tabular}
    \end{center}
    The joint distribution $(\nu_2 \bmod 3,\, \nu_3 \bmod 3)$ has
    $\chi^2 = 0.43$ on 8~degrees of freedom, consistent with perfect
    uniformity.  The singular cells are not attracted to any 3-adic
    sublattice.
\end{empirical}


\subsection{Singular cell chains along the $\log_3 2$ direction}

The spatial structure of singular cells, however, is highly non-trivial.

\begin{empirical}[Connected components]\label{obs:components}
    The 188~singular cells at $k = 81$ form 65~connected components
    (8-connected on the torus):
    \begin{center}
    \renewcommand{\arraystretch}{1.15}
    \begin{tabular}{rc}
        \toprule
        Component size & Count \\
        \midrule
        37 & 1 \\
        18 & 3 \\
        10 & 1 \\
        7 & 4 \\
        4 & 1 \\
        1 (isolated) & 55 \\
        \bottomrule
    \end{tabular}
    \end{center}
    There are no size-2 pairs.  The 133~cells in components of
    size~$\geq 3$ form extended chains.
\end{empirical}

\begin{empirical}[Chain orientation converges to $\log_3 2$]\label{obs:chain-slope}
    Principal component analysis of each chain of size~$\geq 3$ yields a
    consistent principal direction:
    \begin{center}
    \renewcommand{\arraystretch}{1.15}
    \begin{tabular}{rccc}
        \toprule
        Size & Chain slope & $|\mathrm{slope} - \log_3 2|$ & Mean $\Delta(a,b)$ \\
        \midrule
        37 & 0.640 & 0.009 & 0.599 \\
        18 ($\times 3$) & 0.640 & 0.009 & 0.622 \\
        10 & 0.638 & 0.007 & 0.631 \\
        7 ($\times 4$) & 0.640 & 0.009 & 0.581 \\
        \bottomrule
    \end{tabular}
    \end{center}
    All chains are aligned at slope $\approx 0.640$, deviating from
    $\log_3 2 = 1/\log_2 3 \approx 0.6309$ by less than $0.01$.  This
    is neither the unstable foliation slope ($\log_2 3 \approx 1.585$)
    nor the stable foliation slope ($-1/\log_2 3 \approx -0.631$), but
    the \emph{positive reflection} of the stable slope: $+\log_3 2$.
\end{empirical}

\begin{remark}[Diophantine interpretation]
    The best rational approximation of $\log_3 2$ with denominator
    dividing $81 = 3^4$ is $51/81 = 17/27 \approx 0.6296$, with
    error $1.3 \times 10^{-3}$.  The chains trace the direction in
    which $\log_3 2$ is approximated by the $3^4$-lattice.  The chain
    lengths (7, 10, 18, 37) may reflect how far one can walk along the
    $\log_3 2$ direction before the Diophantine approximation error
    accumulates enough to push $p_{\mathrm{odd}}$ back above the
    singular threshold.
\end{remark}


\subsection{Even-dominated dynamics at singular cells}

\begin{empirical}[Singular cells are almost pure-even]\label{obs:singular-podd}
    The parity ratio at singular cells is radically different from the
    global average:
    \begin{center}
    \renewcommand{\arraystretch}{1.15}
    \begin{tabular}{lcc}
        \toprule
        Cell class & Mean $p_{\mathrm{odd}}$ & Std.\ dev. \\
        \midrule
        Global & 0.3323 & --- \\
        Non-singular ($k{=}81$) & 0.3070 & 0.060 \\
        \textbf{Singular} ($k{=}81$) & \textbf{0.0595} & 0.073 \\
        \bottomrule
    \end{tabular}
    \end{center}
    Of the 188~singular cells, 187 (99.5\%) have
    $p_{\mathrm{odd}} < 0.3323$.  These are torus positions where
    trajectories almost exclusively take the even branch (halving),
    rarely encountering an odd step.
\end{empirical}

\begin{keyinsight}[Even tunnels along $\log_3 2$]
    The singular cells at $k = 81$ form coherent chains of
    predominantly even-step dynamics, threaded along the $\log_3 2$
    direction on the torus.  On the pure-$3^4$ lattice, the even
    branch $x \mapsto x/2$ has no resonance with the lattice
    structure ($\gcd(2, 81) = 1$), while the odd branch
    $x \mapsto 3x + 1$ increments $\nu_3$ by~$1$ with perfect
    lattice alignment.  The chains mark exactly where the irrationality
    of $\log_3 2$ creates coherent ``tunnels'' of predominantly even
    dynamics---regions where the 3-adic lattice cannot approximate the
    conjugate winding slope.
\end{keyinsight}


\subsection{Spatial and foliation-distance analysis}

\begin{empirical}[Spatial clustering]\label{obs:spatial-clustering}
    Singular cells exhibit strong spatial clustering on the $k{=}81$ torus:
    \begin{center}
    \renewcommand{\arraystretch}{1.15}
    \begin{tabular}{lc}
        \toprule
        Statistic & Value \\
        \midrule
        Expected neighbours (random) & 0.23 \\
        Observed mean neighbours & 1.45 \\
        Enrichment factor & 6.3$\times$ \\
        \bottomrule
    \end{tabular}
    \end{center}
    The nearest-neighbour count is $6.3\times$ the random expectation.
    Half of all singular cells have exactly 2~neighbours, forming the
    interior of chains.
\end{empirical}

\begin{empirical}[Enrichment near unstable foliation]\label{obs:foliation-enrichment}
    Singular cells concentrate near the \emph{unstable} foliation
    (slope $\log_2 3$):
    \begin{center}
    \renewcommand{\arraystretch}{1.15}
    \begin{tabular}{rcc}
        \toprule
        Distance bin (unstable) & Singular/Total & Enrichment \\
        \midrule
        $[0, 4.3)$ & 60/1313 & 1.59$\times$ \\
        $[4.3, 8.6)$ & 47/1311 & 1.25$\times$ \\
        $[8.6, 13.0)$ & 30/1312 & 0.80$\times$ \\
        $[13.0, 17.3)$ & 26/1312 & 0.69$\times$ \\
        $[17.3, 21.6)$ & 25/1312 & 0.66$\times$ \\
        \bottomrule
    \end{tabular}
    \end{center}
    The enrichment decreases monotonically from $1.59\times$ to
    $0.66\times$ with distance from the unstable foliation.  No
    comparable enrichment is observed with respect to the stable
    foliation (enrichment factors $0.93$--$1.12\times$, essentially
    flat).
\end{empirical}

\begin{remark}[Geometric picture]
    The singular cells are chains \emph{aligned along $\log_3 2$}
    that \emph{concentrate near the $\log_2 3$ foliation}.  These
    two directions are not identical: the chain direction ($32.6°$
    from the $\nu_2$-axis) differs from the unstable foliation
    ($57.8°$) by $25.2°$.  The chains thread along the conjugate
    slope, but preferentially populate the neighbourhood of the
    unstable foliation.
\end{remark}


\subsection{Phase transition at the saturation boundary}

\begin{empirical}[Two regimes of singular cells]\label{obs:two-regimes}
    The character of singular cells reverses sharply at the boundary
    where the branch fraction drops below~$1$:

    \medskip
    \noindent\textbf{Regime~1} ($k \leq 81$, 100\% branch fraction):
    \begin{center}
    \renewcommand{\arraystretch}{1.15}
    \begin{tabular}{rcccccc}
        \toprule
        $k$ & Fact. & Sing. & Max chain & Chain slope & Mean $p_{\mathrm{odd}}$ & \% below global \\
        \midrule
        48 & $2^4 \cdot 3$ & 26 & 7 & 0.673 & 0.170 & 100\% \\
        54 & $2 \cdot 3^3$ & 44 & 10 & 0.732 & 0.124 & 97.7\% \\
        72 & $2^3 \cdot 3^2$ & 95 & 24 & 0.638 & 0.036 & 98.9\% \\
        81 & $3^4$ & 188 & 37 & 0.640 & 0.060 & 99.5\% \\
        \bottomrule
    \end{tabular}
    \end{center}

    \noindent\textbf{Regime~2} ($k \geq 108$, $<100$\% branch fraction):
    \begin{center}
    \renewcommand{\arraystretch}{1.15}
    \begin{tabular}{rcccccc}
        \toprule
        $k$ & Fact. & Sing. & Max chain & Chain slope & Mean $p_{\mathrm{odd}}$ & \% below global \\
        \midrule
        108 & $2^2 \cdot 3^3$ & 65 & 10 & 0.791 & 0.686 & 0.0\% \\
        144 & $2^4 \cdot 3^2$ & 127 & 9 & 0.760 & 0.695 & 0.0\% \\
        \bottomrule
    \end{tabular}
    \end{center}
\end{empirical}

\begin{keyinsight}[Phase transition in singular cell character]
    The two regimes exhibit a complete reversal:
    \begin{enumerate}[label=(\roman*)]
        \item \textbf{Below saturation} ($k \leq 81$): Singular cells
        have very low $p_{\mathrm{odd}}$ ($0.036$--$0.170$), forming
        ``even tunnels'' aligned at slope $\log_3 2$.  The chain slope
        converges to $\log_3 2$ with error $< 0.01$ at $k = 72$ and
        $k = 81$.

        \item \textbf{Above saturation} ($k \geq 108$): Singular cells
        flip to high $p_{\mathrm{odd}}$ ($0.686$--$0.695$), forming
        ``odd outliers'' with shorter chains and slope drifting away from
        $\log_3 2$.

        \item The transition coincides with the appearance of empty cells.
        On a fully saturated torus, the singular set reveals the
        \emph{intrinsic} foliation singularities of the Collatz dynamics.
        On a partially empty torus, the singular set instead reflects the
        \emph{boundary} of the occupied region.
    \end{enumerate}
\end{keyinsight}


\subsection{Diophantine foliation analysis}

The rational approximation quality of $\log_3 2$ by fractions
$p/3^b$ varies dramatically with $b$.  Only certain exponents yield
\emph{record approximations} (strictly improving on all smaller powers
of~3):

\begin{center}
\begin{tabular}{ccccr}
\toprule
$b$ & $3^b$ & Best $p/3^b$ & Approx.\ value & $|p/3^b - \log_3 2|$ \\
\midrule
1  & 3       & $2/3$        & 0.66667 & $3.57 \times 10^{-2}$ \\
3  & 27      & $17/27$      & 0.62963 & $1.30 \times 10^{-3}$ \\
6  & 729     & $460/729$    & 0.63100 & $7.16 \times 10^{-5}$ \\
9  & 19{,}683 & $12419/19683$ & 0.63095 & $2.08 \times 10^{-5}$ \\
11 & 177{,}147 & $111777/177147$ & 0.63093 & $1.76 \times 10^{-6}$ \\
\bottomrule
\end{tabular}
\end{center}

\noindent
The levels $3^4 = 81$ through $3^5 = 243$ all use the same best
approximant $17/27$; the next jump in accuracy occurs at $3^6 = 729$,
which is $18\times$ more precise.  We add $k = 729$ and $k = 19{,}683 = 3^9$
(a further $3.4\times$ improvement, $12419/19683$) as additional
torus levels to probe this Diophantine structure and measure the
scaling of strip width across Diophantine jumps.

\begin{observation}[Frozen branch count]
\label{obs:frozen-branch}
    At $N = 10^9$, the number of branch cells stabilises at exactly
    $13{,}688$ for all $k \geq 144$, including the Diophantine level
    $k = 729$ ($13{,}688$ branch cells of $531{,}441$ total, branch
    fraction $2.58\%$).  The branch locus
    has \emph{finite complexity}: the set of trajectory signatures
    producing branching behaviour is fully captured by $k \approx 144$,
    and larger moduli simply embed the same set in a sparser grid.
\end{observation}

\begin{remark}[Branch count convergence with $N$]
\label{rem:branch-convergence}
    The ``frozen'' property refers to stability across $k$
    (resolution-independence), not convergence in~$N$ (trajectory count).
    As $N$ increases, rare trajectories may visit previously unoccupied
    cells with both parities, converting pure cells to branch cells.
    The refactored code at $N = 10^8$ yields $9{,}415$ branch cells;
    the original run at $N = 10^9$ yielded $13{,}688$.  A definitive
    $N = 10^{10}$ run is in progress to establish convergence.
    The qualitative structure---resolution-independent branch count,
    finite topological complexity---is robust across all $N$ tested.
\end{remark}

\begin{observation}[Foliation depletion at Diophantine levels]
\label{obs:foliation-depletion}
    The enrichment of branch cells on the unstable foliation reverses
    at large~$k$:
    \begin{center}
    \begin{tabular}{crrrrr}
    \toprule
    $k$ & branch & on foliation & branch $\cap$ fol.\ & expected & enrichment \\
    \midrule
    6--72  & all     & ---  & all   & ---   & 1.000 \\
    729    & 13{,}688 & 1{,}367 & 2     & $\sim$35 & 0.057 \\
    \bottomrule
    \end{tabular}
    \end{center}
    At $k \leq 72$ (full saturation), every foliation cell is
    trivially a branch cell.  At $k = 729$, the branch locus is
    $17\times$ \emph{depleted} on the unstable foliation---branch cells
    actively avoid the foliation lines.
\end{observation}

\begin{observation}[Diagonal band structure at $k = 729$]
\label{obs:diagonal-band}
    Despite occupying only $2.8\%$ of the $729 \times 729$ torus, all
    non-empty cells at $k = 729$ lie on a single coherent diagonal band
    of width ${\sim}20$~cells, aligned at slope $\approx \log_3 2$.
    Within this band:
    \begin{itemize}
        \item Branch cells (gold) occupy the core,
        \item Pure-even cells (1{,}269) flank the branch cells on both
              sides, forming ``laminar tunnel walls,''
        \item Pure-odd cells (41) are rare and peripheral.
    \end{itemize}
    The band is \emph{offset} from the exact unstable foliation by
    $15$--$30$~cell units in perpendicular distance.  The radial
    density profile from the foliation shows non-empty cells peaking
    at distance ${\sim}20$, not at distance~$0$.
\end{observation}

\begin{observation}[Foliation shadow]
\label{obs:foliation-shadow}
    The offset between the non-empty band and the exact foliation can
    be understood via the Diophantine approximation.  At $k = 729$,
    the best rational approximation $460/729 \approx 0.63100$ exceeds
    $\log_3 2 \approx 0.63093$ by $7.16 \times 10^{-5}$.  Over
    $729$~cells, this error accumulates to
    $729 \times 7.16 \times 10^{-5} \approx 0.052$~wrapping cycles,
    producing a systematic drift of the trajectory band away from the
    irrational foliation.  The band follows the \emph{rational
    approximation} $\nu_3 = (460/729)\,\nu_2$, not the irrational
    foliation $\nu_3 = \log_3 2 \cdot \nu_2$.  We call this the
    \emph{foliation shadow}: trajectories are enslaved to the best
    rational approximation, creating a narrow corridor that shadows but
    does not coincide with the true foliation.
\end{observation}

\begin{observation}[Parity distribution at Diophantine levels]
\label{obs:dioph-parity}
    The distribution of $p_{\mathrm{odd}}$ among branch cells at
    $k = 729$ is heavily left-skewed relative to $k = 81$, with the
    bulk of the mass below $p_{\mathrm{odd}} = 0.1$.  This confirms
    that the Diophantine branch cells are even-dominated: the better
    the rational approximation, the more the dynamics within the
    foliation shadow are biased toward even (halving) steps.
\end{observation}

\begin{keyinsight}[Diophantine quality controls foliation geometry]
    The quality of the rational approximation $p/3^b \approx \log_3 2$
    determines three geometric features of the branch locus on the
    $3^b$-torus:
    \begin{enumerate}[label=(\roman*)]
        \item \textbf{Band width:} Better approximations produce
        tighter trajectory confinement.  At $k = 81$ (error
        $1.3 \times 10^{-3}$), all cells are visited; at $k = 729$
        (error $7.2 \times 10^{-5}$), only $2.8\%$ of cells are
        non-empty.
        \item \textbf{Foliation offset:} The trajectory band follows
        the rational approximation, not the irrational foliation,
        with the offset proportional to the approximation error.
        \item \textbf{Parity lamination:} Within the band, even-step
        cells dominate, with branch cells (mixed parity) confined to
        the core and pure-even cells forming the walls---a laminar
        structure invisible at saturated levels.
    \end{enumerate}
\end{keyinsight}


\subsection{Parity transition tracking and ``11'' verification}\label{ssec:transitions}

To directly verify the SFT forbidden bigram ``11''
(Section~\ref{sec:parity-sft}) at the level of individual torus cells,
and to compute per-cell transition valence for the Euler characteristic
programme, we extend the branch locus computation with per-cell parity
transition grids.

\paragraph{Setup.}
For each target level $k \in \{81, 108, 144, 729\}$, we track three
additional $k \times k$ grids recording the number of
\emph{even$\to$even}, \emph{even$\to$odd}, and \emph{odd$\to$even}
transitions at each cell $(r_2, r_3)$.  A transition is recorded at
the cell of the current step, with the transition type determined by
the previous and current parities.  No \emph{odd$\to$odd} grid is
needed: after an odd Collatz step $x \mapsto 3x+1$, the result is
always even, so the next step is necessarily even.  Memory cost:
$3 \times k^2 \times 8$~bytes per target level, ${\sim}52$~MB total.

\begin{observation}[``11'' forbidden bigram verified on torus cells]
\label{obs:oo-verified}
    At $N = 10^8$ ($17.9 \times 10^9$ total steps), the transition
    counts at all four target levels satisfy
    \[
        \mathrm{ee} + \mathrm{eo} + \mathrm{oe}
        = N_{\mathrm{steps}} - N_{\mathrm{trajs}}
        = 17{,}823{,}493{,}584
    \]
    exactly, with zero residual.  Since the only missing transition
    type is odd$\to$odd, this confirms $\mathrm{oo} = 0$ at every
    cell---the ``11'' forbidden bigram of the parity SFT
    (Section~\ref{sec:parity-sft}) is verified cell-by-cell on the
    torus.
\end{observation}

\begin{observation}[Transition fractions and valence]
\label{obs:trans-fractions}
    The global transition fractions are
    $\mathrm{ee} : \mathrm{eo} : \mathrm{oe} = 0.335 : 0.331 : 0.334$,
    nearly equal thirds.  The slight asymmetry reflects the departure
    of $p_{\mathrm{odd}} \approx 0.332$ from the equilibrium value
    $1/(1 + \log_2 3) \approx 0.387$ at finite~$N$.

    The \emph{transition valence} of a cell---the number of distinct
    transition types (ee, eo, oe) observed---ranges from 0 (empty) to 3
    (mixed).  At $k = 729$, the valence map reveals internal structure
    within the foliation shadow: cells in the strip core have valence~3,
    while cells at the strip edges have valence 1 or~2, with the pattern
    modulated by the $460/729$ rational lattice.
\end{observation}

\begin{observation}[``11'' successor analysis]
\label{obs:successor}
    For each odd-visited cell $(r_2, r_3)$ at $k = 729$, we classify
    the successor cell $(r_2, (r_3 + 1) \bmod k)$---the cell where
    the trajectory lands after the odd step.  At $N = 10^8$:
    \begin{center}
    \begin{tabular}{lrr}
    \toprule
    Successor type & Count & Fraction \\
    \midrule
    branch (both parities) & 8{,}832 & 91.2\% \\
    pure\_even              & 856    & 8.8\% \\
    pure\_odd               & 0      & 0\% \\
    empty                  & 0      & 0\% \\
    \bottomrule
    \end{tabular}
    \end{center}
    The successor of an odd visit is \emph{never} a pure-odd cell,
    confirming that the ``11'' constraint maps directly onto the
    pure-even tunnel walls of the foliation shadow.  The 8.8\% of
    successors landing in pure-even cells correspond to the strip
    boundary, where the odd step shifts $r_3$ by~1 and pushes the
    trajectory outside the branch core.  This boundary structure
    manifests as ``blue tendrils'' extending from the branch strip
    in the direction of the stable foliation.
\end{observation}

\begin{observation}[Shadow offset spatial structure]
\label{obs:shadow-spatial}
    The shadow offset $\delta = d_{\mathrm{irr}} - d_{\mathrm{rat}}$
    (distance from irrational foliation minus distance from rational
    foliation) exhibits a crisp sign change across the foliation strip
    at $k = 729$:
    \begin{itemize}
        \item Cells \emph{above} the irrational foliation (larger $\nu_2$
              for given $\nu_3$): $\delta < 0$ (closer to irrational).
        \item Cells \emph{below} the irrational foliation: $\delta > 0$
              (closer to rational).
    \end{itemize}
    The boundary between positive and negative offset \emph{is} the
    irrational foliation line itself.  The offset magnitude grows
    linearly along the strip as
    $\Delta s \cdot r_2 / \sqrt{1 + s^2}$,
    where $\Delta s = 460/729 - \log_3 2 \approx 7.16 \times 10^{-5}$,
    reaching a maximum of ${\approx}\,0.044$~cell units at $r_2 = 729$.
    Since the rational slope exceeds the irrational by $\Delta s > 0$,
    the rational foliation sits slightly below the irrational in
    the $(\nu_3, \nu_2)$ plot, producing the asymmetry: cells above
    the irrational line are farther from the rational line (and vice
    versa).  The two foliations are at most $0.044$~cells apart, so
    essentially no branch cells sit between them---every cell is
    unambiguously on one side.
\end{observation}

\begin{keyinsight}[SFT--tunnel correspondence at $k = 729$]
    The per-cell transition analysis establishes a direct link between
    the symbolic dynamics (parity SFT, Section~\ref{sec:parity-sft})
    and the geometric structure (foliation shadow,
    Observation~\ref{obs:foliation-shadow}):
    \begin{enumerate}[label=(\roman*)]
        \item The forbidden bigram ``11'' is verified at every cell
              of every target level (implied $\mathrm{oo} = 0$).
        \item The successors of odd-visited cells are never pure-odd,
              confirming that the ``11'' constraint creates the
              pure-even tunnel walls flanking the branch core.
        \item The shadow offset cleanly bisects the strip into cells
              closer to the rational vs.\ irrational foliation,
              with the irrational line as the precise dividing surface.
    \end{enumerate}
\end{keyinsight}


\subsection{Synthesis}

\begin{mainresult}[Branch locus structure on finite tori]
\label{res:branch-locus}
    The per-step Collatz dynamics on $(\Z/k\Z)^2$ for $k = 2^a \cdot 3^b$
    ($0 \leq a,b \leq 4$) and the Diophantine level $k = 3^6 = 729$,
    computed from $10^8$--$10^9$ trajectories with per-cell parity
    transition tracking, reveals:
    \begin{enumerate}[label=(\roman*)]
        \item \textbf{Universal branching at saturation.}  All levels
        $k \leq 81$ have 100\% branch cells: every torus position is visited
        by both even and odd steps.  Above $k = 108$, the branch count
        stabilises at exactly $13{,}688$---the branch locus has finite
        complexity.

        \item \textbf{Singular chains along $\log_3 2$.}  At the fully
        saturated levels $k = 72$ and $k = 81$, singular cells form
        extended chains (up to 37~cells) aligned at slope
        $\log_3 2 \approx 0.631$ with precision $< 0.01$.  These chains
        mark coherent regions of predominantly even dynamics.

        \item \textbf{Foliation enrichment and depletion.}  At saturated
        levels ($k \leq 72$), singular cells are $1.6\times$ enriched
        near the unstable foliation.  At the Diophantine level $k = 729$,
        branch cells are $17\times$ \emph{depleted} on the foliation
        (enrichment $0.057$)---a complete reversal.

        \item \textbf{Foliation shadow.}  At $k = 729$, non-empty cells
        form a single diagonal band of width ${\sim}20$~cells at slope
        $\log_3 2$, offset by $15$--$30$~cell units from the exact
        foliation.  The band follows the best rational approximation
        $460/729 \approx 0.63100$, not the irrational $\log_3 2$.

        \item \textbf{Parity lamination.}  Within the foliation shadow,
        even-step cells dominate: branch cells (mixed parity) form the
        core, flanked by pure-even walls.  The $p_{\mathrm{odd}}$
        distribution at $k = 729$ is heavily left-skewed ($< 0.1$).

        \item \textbf{Phase transition.}  The singular cell character
        reverses at the saturation boundary: even-dominated tunnels
        ($p_{\mathrm{odd}} \approx 0.06$) below, odd-dominated outliers
        ($p_{\mathrm{odd}} \approx 0.69$) above.

        \item \textbf{No sublattice structure.}  Despite $k = 81 = 3^4$
        being a pure prime power, the singular cells are uniformly
        distributed across all 3-adic sublattices.

        \item \textbf{``11'' verification.}  Per-cell transition grids at
        target levels $k \in \{81, 108, 144, 729\}$ confirm that
        $\mathrm{oo} = 0$ at every cell ($17.8 \times 10^9$ transitions
        verified).  The successors of odd-visited cells at $k = 729$ are
        never pure-odd (91.2\% branch, 8.8\% pure-even), directly linking
        the parity SFT forbidden bigram to the pure-even tunnel walls of
        the foliation shadow.

        \item \textbf{Shadow offset bisection.}  The shadow offset
        $d_{\mathrm{irr}} - d_{\mathrm{rat}}$ exhibits a crisp sign
        change at the irrational foliation, bisecting the strip into cells
        closer to the rational ($460/729$) vs.\ irrational ($\log_3 2$)
        foliation.  The offset grows linearly as
        $\Delta s \cdot r_2 / \sqrt{1 + s^2}$ with maximum $0.044$~cell
        units, confirming that the two foliations are nearly coincident at
        $k = 729$.
    \end{enumerate}
\end{mainresult}

\begin{remark}[Irrational foliation passes through the branch strip]
\label{rem:foliation-through-strip}
    A natural question is whether the irrational foliation (the
    equilibrium line $\nu_3 = \log_3 2 \cdot \nu_2$) is geometrically
    separated from the branch cells.  Direct computation at $k = 729$
    shows it is \emph{not}: $180$ branch cells lie within $0.5$~cell
    units of the irrational foliation, with the minimum distance
    $\approx 0$.  The irrational line passes through the densest part
    of the branch strip.

    The ``15--30 cell offset'' (Observation~\ref{obs:foliation-depletion})
    refers to the peak of the \emph{radial density profile} measured
    from the \emph{unstable} foliation (slope $\log_2 3$), not the
    distance from the trajectory band to the irrational line itself.
    This is consistent with the shadow offset bisection
    (Observation~\ref{obs:shadow-spatial}), which shows the irrational
    foliation as the \emph{dividing surface within} the strip.

    The mechanism by which trajectories are forced to descend is
    therefore not geometric exclusion ($A \cap L = \emptyset$, where
    $A$ is the branch locus and $L$ the equilibrium line).
    Instead, the pure-even tunnel walls flanking the branch core act
    as \textbf{reflecting boundaries}: when a trajectory approaches
    the tunnel edge, it encounters pure-even cells that force $n/2$
    divisions, pushing the walk $u(t)$ downward.  The Baker bound
    guarantees that these walls persist at every scale---the tunnel
    has minimum width bounded away from zero---ensuring the reflecting
    boundary is permanent.
\end{remark}

\begin{remark}[The role of $\log_3 2$ and its Diophantine approximants]
    The appearance of $\log_3 2 = 1/\log_2 3$ as the chain direction
    is significant.  The two fundamental slopes of the Collatz foliation
    are $\log_2 3$ (unstable) and $-1/\log_2 3$ (stable).  The singular
    chains are aligned along the \emph{magnitude} of the stable slope,
    but with positive sign---neither foliation direction, but the
    conjugate winding ratio $\log_3 2$ that measures the average number
    of halving steps per tripling step.  This is the same quantity that
    governs the drift along the unstable eigendirection
    (Section~\ref{sec:anosov}).

    The Diophantine analysis (Section~\ref{sec:branch-locus}.\ref{obs:foliation-shadow})
    confirms this interpretation: at levels where $p/3^b$ is a record
    approximation to $\log_3 2$, trajectories are enslaved to the
    rational foliation $\nu_3 = (p/3^b)\,\nu_2$, forming a narrow
    band that shadows the irrational line.  The quality of the
    approximation controls the band width (tighter $\Rightarrow$ narrower)
    and the foliation offset (tighter $\Rightarrow$ closer, but never
    coincident).  The sequence of Diophantine best approximants
    $b = 1, 3, 6, 9, 11, \ldots$ defines a hierarchy of torus
    levels at which the foliation shadow is maximally constrained,
    potentially connecting to the continued-fraction expansion of
    $\log_3 2$ and the Baker-type lower bounds on $|2^m - 3^n|$.
\end{remark}

\begin{remark}[Implementation]
    The branch locus computation is implemented in
    \texttt{branch\_locus.c}, which performs per-step streaming
    accumulation across 26~torus levels simultaneously (25~standard
    levels $2^a \cdot 3^b$ with $0 \leq a,b \leq 4$, plus the
    Diophantine level $3^6 = 729$).  Per-cell parity transition grids
    (ee, eo, oe) are tracked at four target levels
    $k \in \{81, 108, 144, 729\}$, adding ${\sim}52$~MB.  Total memory:
    ${\sim}100$~MB.  At $N = 10^8$, the program processes
    $1.79 \times 10^{10}$~steps in 9~minutes
    (${\sim}34 \times 10^6$~steps/s), writing per-cell data to CSV
    (\texttt{branch\_cells.csv}, \texttt{branch\_transitions.csv},
    \texttt{branch\_shadow.csv}) for offline analysis.  Companion plots
    are generated by \texttt{plot\_branch\_locus.py} (standard levels)
    and \texttt{plot\_diophantine.py} (Diophantine and transition
    analysis, 5~figures).  A definitive $N = 10^{10}$ run is in progress.
\end{remark}


%%%%%%%%%%%%%%%%%%%%%%%%%%%%%%%%%%%%%%%%%%%%%%%%%%%%%%%%%%%%%%%%%%%%%%%%%%%%%%%
% References
%%%%%%%%%%%%%%%%%%%%%%%%%%%%%%%%%%%%%%%%%%%%%%%%%%%%%%%%%%%%%%%%%%%%%%%%%%%%%%%

\begin{thebibliography}{99}

\bibitem{Lagarias1985}
J.~C.~Lagarias,
``The $3x+1$ problem and its generalizations,''
\emph{American Mathematical Monthly} \textbf{92} (1985), no.~1, 3--23.

\bibitem{Lagarias2010}
J.~C.~Lagarias (ed.),
\emph{The Ultimate Challenge: The $3x+1$ Problem},
American Mathematical Society, Providence, RI, 2010.

\bibitem{Terras1976}
R.~Terras,
``A stopping time problem on the positive integers,''
\emph{Acta Arithmetica} \textbf{30} (1976), no.~3, 241--252.

\bibitem{Wirsching1998}
G.~J.~Wirsching,
\emph{The Dynamical System Generated by the $3n+1$ Function},
Lecture Notes in Mathematics, vol.~1681, Springer-Verlag, Berlin, 1998.

\bibitem{Bohm1978}
C.~B\"ohm and G.~Sontacchi,
``On the existence of cycles of given length in integer sequences like $x_{n+1} = x_n/2$ if $x_n$ even, and $x_{n+1} = 3x_n + 1$ otherwise,''
\emph{Atti Accad.\ Naz.\ Lincei Rend.\ Cl.\ Sci.\ Fis.\ Mat.\ Natur.} \textbf{64} (1978), no.~3, 260--264.

\bibitem{Everett1977}
C.~J.~Everett,
``Iteration of the number-theoretic function $f(2n) = n$, $f(2n+1) = 3n+2$,''
\emph{Advances in Mathematics} \textbf{25} (1977), no.~1, 42--45.

\bibitem{Tao2022}
T.~Tao,
``Almost all orbits of the Collatz map attain almost bounded values,''
\emph{Forum of Mathematics, Pi} \textbf{10} (2022), e12.

\bibitem{Thurston1988}
W.~P.~Thurston,
``On the geometry and dynamics of diffeomorphisms of surfaces,''
\emph{Bulletin of the American Mathematical Society} \textbf{19} (1988), no.~2, 417--431.

\bibitem{Oliveira2015}
T.~Oliveira e~Silva,
``Empirical verification of the $3x+1$ and related conjectures,''
in~\cite{Lagarias2010}, pp.~189--207.

\bibitem{Vietoris1927}
L.~Vietoris,
``\"Uber den h\"oheren Zusammenhang kompakter R\"aume und eine Klasse von zusammenhangstreuen Abbildungen,''
\emph{Mathematische Annalen} \textbf{97} (1927), 454--472.

\bibitem{LindMarcus2021}
D.~Lind and B.~Marcus,
\emph{An Introduction to Symbolic Dynamics and Coding},
2nd ed., Cambridge University Press, Cambridge, 2021.

\end{thebibliography}

\end{document}
